\documentclass[11pt]{article}
\usepackage[a4paper,margin=27mm]{geometry}
\usepackage[T1]{fontenc}
\usepackage{lmodern,microtype,amsmath,amssymb,amsthm,mathtools,enumitem}
\usepackage[hidelinks,hypertexnames=false]{hyperref}
\newtheorem{theorem}{Theorem}
\newtheorem{lemma}[theorem]{Lemma}
\newtheorem{proposition}[theorem]{Proposition}
\newtheorem{corollary}[theorem]{Corollary}
\theoremstyle{definition}\newtheorem{definition}[theorem]{Definition}
\theoremstyle{remark}\newtheorem{remark}[theorem]{Remark}
\newcommand{\R}{\mathbb R}
\newcommand{\Dscr}{\mathcal D}
\newcommand{\supp}{\operatorname{supp}}
\newcommand{\debt}{\mathrm{debt}}
\newcommand{\cred}{\mathrm{credit}}
\newcommand{\sil}{\mathrm{silence}}
\newcommand{\Bfam}{\mathcal B}
\newcommand{\Jid}{\mathcal J}
\newcommand{\Iid}{\mathcal I}
\newcommand{\Uf}{\mathcal U}
\newcommand{\UU}{\mathfrak U}
\newcommand{\Alg}{\mathfrak A}
\newcommand{\Lo}[3]{\mathrm L^{#1}_{#2}(#3)}
\newcommand{\Hi}[3]{\mathrm H^{#1}_{#2}(#3)}
\newcommand{\CN}{\mathrm{CN}}
\newcommand{\DI}{\mathrm{DI}}
\newcommand{\pw}{\mathrm{pw}}
\title{Complete Rankings for Any Number of Worlds}
\author{}\date{11 October 2026}
\begin{document}\maketitle

\noindent
This note continues FSE (\emph{Fineness Spreads Along Edges}), CCF
(\emph{Closed Choice Forces Fineness}) and the three-world note
(\emph{Three Complete Rankings Cannot Cover}), and uses, through them,
results of CBOI (\emph{Closed Choice Beyond Ordinary Integrals}). The
three-world note solves problem~P1 of FSE, Section~1.7, for three worlds.
Here we solve it for every finite number of worlds: if $n\ge3$, every
ranking is a complete preorder satisfying (G) and (IB), and the selection
axiom (S) holds, then the vacuous model $V$ is dominated.

\section{Summary and consequences}\label{sec:summary}

\begin{enumerate}[leftmargin=*,itemsep=3pt]
\item \emph{Main theorem} (Theorem~\ref{thm:main}). Let $n\ge3$, let every
ranking be a complete preorder satisfying (G) and (IB), and assume (S).
Then $V$ is dominated, and hence dominated by a selected coherent model.
\item \emph{Clique theorem} (Theorem~\ref{thm:clique},
Corollary~\ref{cor:consequences}). Under the same hypotheses, let $D_K$ be
a coherent model whose credal set $K$ contains at least three vertices of
the simplex. If $D_K$ is undominated, then every edge between two vertices
in $K$ has the endpoint property and is fine, and some edge of the simplex
does not have the endpoint property. Hence, if some edge carries a
selection whose credal set omits both endpoints, every coherent model whose
credal set contains these two endpoints and a third vertex is dominated.
\item \emph{Low-cone heaviness in $n$ worlds} (Theorem~\ref{thm:LH}). Under
the same hypotheses, and with no assumption on $V$: at every world, credit
on any non-null set compensates some non-null debt, and debt on some part
of any non-null set is not compensated by some non-null credit. For $n=3$
this is FSE's Theorem~56.
\item \emph{Own regions and the endpoint dichotomy}
(Theorems~\ref{thm:own} and~\ref{thm:dichotomy}). Let $D_K$ be undominated,
and let $S(K)$, the set of worlds whose vertices lie in $K$, have at least
three members. At every two-sided selection on an edge between two worlds
of $S(K)$, the debt order of each of the two edge worlds is non-fine on its
own odd region, relative to credit on deep tails of the odd region of any
third world of $S(K)$. Either every edge between two worlds of $S(K)$ has
the endpoint property, or none has.
\item \emph{Rich far ends and counting} (Propositions~\ref{prop:rich}
and~\ref{prop:finite}). The two main steps of the three-world note hold in
$n$ worlds, on any three worlds of $S(K)$.
\end{enumerate}

The device that makes the passage to $n$ worlds possible is the family of
regions $O_v$, $v=1,\dots,n$: $O_v$ is the set of mixed problems at which
world~$v$ disagrees with every other world. On $O_v$ the option bad at $v$
is good at every other world. So a change made of such options on $O_u$,
$O_v$ and $O_w$ improves every world outside $\{u,v,w\}$ strictly, and the
dominance arguments of FSE and of the three-world note, which concern
three worlds, can be run on these regions without losing the other worlds.
Comparisons at a single world may use the larger three-world regions
$D_{vw}\cap D_{vu}$; dominance arguments may not (Remark~\ref{rem:regions}).

\emph{Outline of the proof.} Suppose $V$ is undominated. By the endpoint
dichotomy with $K=\Delta$, either every edge has the endpoint property, and
then CCF's Theorem~16 makes $V$ dominated, or no edge has it, and every
edge carries a two-sided selection. In the second case the proof of the
three-world note runs in $n$ worlds on three fixed worlds $1,2,3$, with the
three-world regions $P_x$ replaced by $O_x$. The own-region theorem yields
uncountably many ultrafilters of a certain kind on $O_2$; the
no-saturated-cut theorem and the counting lemma of the three-world note,
which hold for every $n$, show that there are only countably many, once
the far end of the edge $[e_1,e_2]$ is known to be rich. Richness is proved
as in the three-world note, with the own-region theorem and low-cone
heaviness in $n$ worlds as inputs. The same argument, with an undominated
coherent model $D_K$ in place of $V$, proves the clique theorem.

\emph{What remains open.} First, CBOI's form of the conclusion for
$n\ge4$: whether every coherent model whose credal set contains three
vertices is dominated. For $n=3$ the only such model is $V$, so this is the
main theorem. For $n\ge4$, by the clique theorem, an undominated such model
$D_K$ needs the endpoint property on every edge between vertices in $K$ and
its failure on some other edge; so $K$ omits some vertex. Our arguments do
not reach the edges that leave the vertices in $K$: if $e_x\notin K$, then
$D_K$ gives the option bad at $x$ on a non-null part of $O_x$, and on all
of $O_x$ if $K$ lies in the face spanned by the vertices of $\Delta$ that it
contains, while the dominance arguments need a base model that is silent
on the regions they change; and CCF's Theorem~16 needs the endpoint
property on every edge (Remark~\ref{rem:open}). Second, incomplete
rankings. For every $n\ge3$, CCF's Lemma~1 makes $V$ dominated, with no
completeness and no (IB) at any world, when every world carries a
one-sided weak extension of CBOI's form, that is, a loss of CBOI's form
satisfying CCF's one-sided bound (W$_k$) (CCF, Section~1); for $n=3$, TAW
does so when two worlds carry such extensions and the third satisfies
(IB). Beyond these cases we know
neither a proof that $V$ is dominated nor an example in which it is
undominated when some ranking is incomplete.

The results of Sections~\ref{sec:LH}--\ref{sec:main} are taken from a
working note that was checked by an independent referee.
Section~\ref{sec:rich} instead adapts the proof of the three-world note's
Proposition~6, which was refereed for $n=3$; in $n$ worlds only its last
step changes. The note as written, including that step, was checked by a
further independent referee. The corrections of both referees are
incorporated. Results that are not in the citable notes are proved in
full. Both referees also checked the sign and region bookkeeping of the
constructions, at every world, by sampling random problems for
$n=3,\dots,7$: the first those of the working note, the second Step~5 of
Proposition~\ref{prop:rich} and the instance of Theorem~\ref{thm:own}
that it uses. A collaborator's note, \emph{Fineness for Any Finite Number of Worlds}, which has not been
refereed, states several of FSE's results for every $n\ge3$, among them
low-cone heaviness, and observes that changes made of the options
$\omega_v$ on the regions $O_v$ (Section~\ref{sec:regions}) improve every
world outside the active triple. We do not use its results.
Theorem~\ref{thm:LH} proves its theorem on low-cone heaviness in full, and
its theorem that one fine edge forces $V$ to be dominated follows from
Proposition~\ref{prop:partner}, Theorem~\ref{thm:dichotomy} and CCF's
Theorem~16.

\section{Setting and inputs}\label{sec:setting}

\subsection{Problems, models and axioms}\label{sec:axioms}

We use the setting and notation of FSE, Section~1, and of the three-world
note, Section~2, as far as they concern an arbitrary number $n\ge3$ of
worlds, and recall what is needed. $\Delta\coloneqq\{p\in\R^n:p\ge0,\
\sum_ip_i=1\}$ has vertices $e_i$. Binary problems $\{g,-g\}$ are
represented by $g\in H\coloneqq\{g\in\R^n:g_1>0\}$, with normalized
Gaussian measure~$\mu$; coordinate hyperplanes are null, and $H'$ is the
set of mixed problems in $H$. Models are measurable and identified modulo
null sets. A model is \emph{regular} if it is never contradictory and
correct on unanimous problems. For a nonempty compact convex
$K\subseteq\Delta$, the coherent model $D_K$ recommends, on a mixed problem
$\{g,-g\}$, the option $h\in\{g,-g\}$ with $\min_{p\in K}p\cdot h>0$ if
there is one, and is silent otherwise (CBOI, Section~1); $K$ is its
\emph{credal set}, and we put $S(K)\coloneqq\{x:e_x\in K\}$. Coherent
models are regular; $V\coloneqq D_\Delta$, and $D_i$ is the ideal model of
world~$i$, which recommends the option good at~$i$ on every problem. All
models below are regular.

On a mixed problem a regular model gives the option good at world~$v$
(\emph{credit} at $v$, rank~$0$), silence (rank~$1$), or the option bad at
$v$ (\emph{debt} at $v$, rank~$2$). A regular model $M$ is determined by
its world-$v$ rank function $r_v(\cdot,M)\colon H'\to\{0,1,2\}$, and every
measurable function $H'\to\{0,1,2\}$ is the world-$v$ rank function of a
regular model. For measurable $Q\subseteq H'$, the models $M+\debt_v(Q)$,
$M+\cred_v(Q)$ and $M+\sil(Q)$ agree with $M$ off $Q$ and have world-$v$
rank $2$, $0$, $1$ on $Q$; $\debt_v(M)$, $\cred_v(M)$ and $\sil(M)$ denote
the sets where $M$ has world-$v$ rank $2$, $0$, $1$. $V$ is silent on $H'$,
and $D_v$ has world-$v$ rank $0$ on $H'$. We write $X\le^v_\pw Y$ if
$r_v(\cdot,X)\le r_v(\cdot,Y)$ on $H'$.

At each world, $\preceq_i$ is a preorder on models; smaller is better, and
$X\prec_iY$ means that $X\preceq_iY$ and not $Y\preceq_iX$. The ranking of
world~$i$ is \emph{complete} if any two models are comparable. (G) at
world~$i$ is guidance monotonicity: $X\le^i_\pw Y$ implies $X\preceq_iY$
(\emph{weak (G)}), and if moreover $r_i(\cdot,X)<r_i(\cdot,Y)$ on a
non-null set, then $X\prec_iY$ (\emph{strict (G)}). (IB) at world~$i$ is
independence under replacement of shared advice; we use it in the form of
the three-world note, Section~2.1:
\begin{enumerate}[label=(IB$'$),leftmargin=*]
\item If $X,Y,X',Y'$ are regular, $\Theta\coloneqq\{X\ne Y\}=\{X'\ne
Y'\}$, and $X=X'$ and $Y=Y'$ on $\Theta$, then $X\preceq_iY\iff
X'\preceq_iY'$ and $X\prec_iY\iff X'\prec_iY'$.
\end{enumerate}
So whether $X\preceq_vY$ depends only on the set $\{X\ne Y\}$ and on the
world-$v$ ranks of $X$ and $Y$ there.

(S) supplies coherent undominated models $C_p$, $p\in\Delta$, competitive
on $\supp p$, with closed sections $A_i(X)\coloneqq\{p:C_p\preceq_iX\}$. A
model $C$ is \emph{competitive} on a set $\mathcal W$ of worlds if for
every model $X$ some $i\in\mathcal W$ has $C\preceq_iX$. A model $M$ is
\emph{dominated} if some model $X$ has $X\preceq_iM$ at every world and
$X\prec_iM$ at some world. By CBOI's Lemma~1, every model has a selected
coherent weak dominator at all worlds; so a dominated model is dominated
by a selected coherent model. By the face lemma (FSE, Section~1.1) the
credal set of $C_p$ lies in the face spanned by $\supp p$; in particular
$C_{e_i}=D_i$. For $q$ on an edge $[e_v,e_w]$, $C_q$ is a \emph{selection
on the edge}; if $q$ lies in the open edge $(e_v,e_w)$, $C_q$ is
competitive on $\{v,w\}$. We use two consequences of (S) (FSE,
Section~1.1; CCF, Section~3; the three-world note, Section~2.1).
\begin{enumerate}[label=(T\arabic*),leftmargin=*,itemsep=2pt]
\item If world $k$ is complete, $Y$ is regular and $Y\prec_kC_{p^*}$, then
$Y\prec_kC_p$ for all $p$ near $p^*$.
\item If $p^*\in[e_v,e_w]$, $X$ is regular and $X\prec_wC_{p^*}$, then
$C_p\preceq_vX$ for all $p$ in the open edge $(e_v,e_w)$ near $p^*$.
\end{enumerate}
(T1) holds because $A_k(Y)$ is closed and does not contain $p^*$, and
completeness turns $C_p\not\preceq_kY$ into $Y\prec_kC_p$. (T2) holds
because $A_w(X)$ is closed and does not contain $p^*$, and $C_p$ is
competitive on $\{v,w\}$.

\subsection{Debt, credit, pairs and fineness}\label{sec:pairs}

For measurable $B,C\subseteq H'$ (FSE, Section~1.2),
\begin{gather*}
 B\preceq^v_dC\ :\Longleftrightarrow\ V+\debt_v(B)\preceq_vV+\debt_v(C),\\
 B\preceq^v_cC\ :\Longleftrightarrow\ V+\cred_v(C)\preceq_vV+\cred_v(B).
\end{gather*}
So $B\preceq^v_dC$ says that debt on $B$ is no worse than debt on $C$, and
$B\preceq^v_cC$ says that credit on $B$ is worth no more than credit on
$C$. Their strict parts are $\prec^v_d$ and $\prec^v_c$; both relations
are transitive. By (IB$'$), $V$ may be replaced by any regular model silent
on $B\cup C$, and $B$ and $C$ may overlap. By weak (G), $B\subseteq C$
implies $B\preceq^v_dC$ and $B\preceq^v_cC$.

A \emph{pair} $(A,T)$ at world $v$ consists of disjoint measurable sets
$A,T\subseteq H'$ (debt part and credit part). It is \emph{SI} (a strict
improvement) if $V+\debt_v(A)+\cred_v(T)\prec_vV$, \emph{weakly SI} if
$V+\debt_v(A)+\cred_v(T)\preceq_vV$, \emph{AC} if it is not SI, and
\emph{strictly AC} if $V\prec_vV+\debt_v(A)+\cred_v(T)$. If world $v$ is
complete, a pair is AC if and only if $V\preceq_vV+\debt_v(A)+\cred_v(T)$,
and a pair that is not weakly SI is strictly AC. By (IB$'$), $V$ may be
replaced by any regular model silent on $A\cup T$. By strict (G), a pair
$(A,T)$ with $A$ null and $T$ non-null is SI, and one with $A$ non-null and
$T$ null is strictly AC. No pair is both SI and strictly AC.

Let $\Alg$ be the measure algebra of $H'$. Ideals and ultrafilters are
those of $\Alg$; so ``$B$ is contained in $C$'' is read modulo null sets
when membership in an ideal is concerned. An ultrafilter \emph{avoids} an
ideal if they are disjoint. For non-null $A\subseteq H'$, $\Jid^v_A$ is
the family of measurable sets contained in a finite union of measurable
sets each $\preceq^v_dA$; for non-null $T$, $\mathcal K^v_T$ is defined
likewise with $\preceq^v_c$. Both are ideals (FSE, Definition~1 and
Lemma~2). For world~$1$ we write $\preceq_d$, $\prec_d$, $\preceq_c$ and
$\Jid_A$. For a non-null measurable $R\subseteq H'$, world $v$'s debt order
is \emph{fine on $R$} if $R\in\Jid^v_A$ for every non-null $A\subseteq R$,
its credit order is fine on $R$ if $R\in\mathcal K^v_T$ for every non-null
$T\subseteq R$, and world $v$ is \emph{fine on $R$} if both are (FSE,
Definition~1). An edge $[e_v,e_w]$ is \emph{fine} if worlds $v$ and $w$
are fine on $D_{vw}$ (FSE, Section~4).

\subsection{Edges and frames}\label{sec:edges}

For distinct worlds $v,w$ put $D_{vw}\coloneqq\{g\in H':g_vg_w<0\}$ and
$t^v_w\coloneqq|g_v|/|g_w|$ on $D_{vw}$, so $t^w_v=1/t^v_w$; the law of
$t^v_w$ under $\mu|_{D_{vw}}$ has a positive density on $(0,\infty)$, and
its level sets are null. A \emph{band} of the edge $[e_v,e_w]$ is a set
$D_{vw}\cap\{s<t^v_w<s'\}$ with $0\le s<s'\le\infty$. For $0<a,b<\infty$,
\[
 \Lo vwa\coloneqq D_{vw}\cap\{t^v_w<a\},\qquad
 \Hi vwb\coloneqq D_{vw}\cap\{t^v_w>b\}
\]
are the \emph{low cone} and the \emph{high tail} of world $v$ on the edge
$[e_v,e_w]$; $\Lo vwa=\Hi wv{1/a}$. On $D_{vw}$ the option good at $v$ is
the option bad at $w$, so the world-$v$ and world-$w$ ranks of every
regular model add up to $2$ (FSE, equation~(1)):
\begin{equation}
 \cred_v=\debt_w\quad\text{and}\quad\debt_v=\cred_w\qquad\text{on }D_{vw};
 \label{eq:swap}
\end{equation}
off $D_{vw}$ the two worlds agree, so their ranks are equal.

\emph{Edge models} (FSE, Section~1.5). By the face lemma a selection on
$[e_v,e_w]$ has a credal interval on that edge, and there are
$0\le\beta\le\alpha\le\infty$ such that it gives credit at $v$ on
$D_{vw}\cap\{t^v_w>\alpha\}$, silence on $D_{vw}\cap\{\beta<t^v_w<\alpha\}$,
debt at $v$ on $D_{vw}\cap\{t^v_w<\beta\}$, and the option good at $v$ and
$w$ off $D_{vw}$. It is written $(\alpha,\beta)$, or
$(\alpha^{(v)},\beta^{(v)})$, \emph{in world $v$'s frame}; in world $w$'s
frame it is $(\alpha^{(w)},\beta^{(w)})=(1/\beta^{(v)},1/\alpha^{(v)})$, by
\eqref{eq:swap}. Regular models of this form are \emph{edge models}. The
model $(0,0)$ is $D_v$ and $(\infty,\infty)$ is $D_w$; since
$\beta\le\alpha$, an edge model with $\alpha=0$ is $D_v$, and one with
$\beta=\infty$ is $D_w$. The credal set of a selection on $[e_v,e_w]$
contains $e_v$ if and only if $\beta=0$, and $e_w$ if and only if
$\alpha=\infty$. A selection on $[e_v,e_w]$ is \emph{two-sided} if
$0<\beta\le\alpha<\infty$. A selection whose credal set contains neither
endpoint lies in the open edge, by the face lemma, and is two-sided (FSE,
Section~1.5). In world $v$'s frame on $[e_v,e_w]$ we write
$E_w(b)\coloneqq(\infty,b)$ for $0\le b\le\infty$: it gives debt at $v$ on
$D_{vw}\cap\{t^v_w<b\}$, silence on the rest of $D_{vw}$, and the option
good at $v$ and $w$ off $D_{vw}$; $E_w(\infty)=D_w$.

The edge $[e_v,e_w]$ has the \emph{endpoint property} (EP$_{vw}$) if every
selection on it has $e_v$ or $e_w$ in its credal set (CCF, Definition~2).
Equivalently, by the last paragraph, no selection on the open edge
$(e_v,e_w)$ is two-sided.

Statements made in FSE or CCF for world~$1$ and an edge $[e_1,e_j]$ hold
for every world $v$ and edge $[e_v,e_w]$, with $D_{vw}$ and $t^v_w$ in
place of $D_{1j}$ and $t_j$, since the axioms are symmetric in the labels
and every object involved is a set of problems; the change of representing
half-space is harmless (FSE, Section~1.3). We use this relabelling without
further comment.

\subsection{The regions \texorpdfstring{$O_v$}{O\_v}}\label{sec:regions}

For each world $v$ put
\[
 O_v\coloneqq\{g\in H':g_vg_x<0\text{ for every }x\ne v\},
\]
the set of mixed problems at which world $v$ disagrees with every other
world, and let $\omega_v$ denote, on $O_v$, the option bad at~$v$. For
measurable $Q\subseteq O_v$, $M+\omega_v(Q)$ is the model that agrees with
$M$ off $Q$ and gives $\omega_v$ on $Q$. For $n=3$, $O_v$ is the region
$P_v$ of FSE, Section~1.4, and $\omega_v$ is the option $o_{yz}$ good
exactly at the two other worlds $y,z$.

\begin{lemma}\label{lem:regions}
Let $n\ge3$.
\begin{enumerate}[label=(\alph*),leftmargin=*,itemsep=2pt]
\item Each $O_v$ is a nonempty open subset of $\R^n$, and the sets
$O_1,\dots,O_n$ are pairwise disjoint. $O_v\subseteq D_{vx}$ for every
$x\ne v$, and $O_v\cap D_{xy}=\varnothing$ for all distinct $x,y\ne v$.
\item On $O_v$, $\omega_v$ is debt at $v$ and credit at every world $x\ne
v$. So for measurable $Q\subseteq O_v$ and $x\ne v$,
$M+\omega_v(Q)=M+\debt_v(Q)=M+\cred_x(Q)$.
\item For distinct $j,k\ne v$, the joint law of $(t^v_j,t^v_k)$ under
$\mu|_{O_v}$ has a strictly positive density on $(0,\infty)^2$, and the law
of $t^v_j$ has a strictly positive density on $(0,\infty)$. Hence every
set $O_v\cap\{(t^v_j,t^v_k)\in G\}$ with $G\subseteq(0,\infty)^2$ Borel of
positive Lebesgue measure is non-null, and so is every nonempty open subset
of $H'$.
\item Let $E$ be a selection on $[e_u,e_v]$ with thresholds $(\alpha,\beta)$
in world $u$'s frame. Then $E$ gives $\omega_v$ on $O_v\cap\{t^u_v>\alpha\}$
and $\omega_u$ on $O_u\cap\{t^u_v<\beta\}$.
\item Let $\mathcal X$ be a set of worlds, and for each $y\in\mathcal X$
let $Q_y\subseteq O_y$ be measurable. If $M$ is silent on
$\bigcup_yQ_y$, then $M+\sum_y\omega_y(Q_y)\le^x_\pw M$ for every world
$x\notin\mathcal X$, with strict inequality of ranks on $\bigcup_yQ_y$; so,
under (G) at $x$, $M+\sum_y\omega_y(Q_y)\prec_xM$ if $\bigcup_yQ_y$ is
non-null.
\item If $x\ne y$ and $e_x,e_y\in K$, then $D_K$ is silent on $D_{xy}$. In
particular, if $|S(K)|\ge2$, then $D_K$ is silent on $O_x$ for every $x\in
S(K)$.
\end{enumerate}
\end{lemma}
\begin{proof}
(a) For $v=1$, $O_v=\{g_1>0,\ g_x<0\text{ for }x\ne1\}$; for $v\ne1$,
$O_v=\{g_v<0,\ g_x>0\text{ for }x\ne v\}$. Both are nonempty open orthants
in $H$. On $O_v$ the worlds $x\ne v$, of which there are at least two,
share one sign, and $v$ has the other; so $v$ is the only coordinate whose
sign differs from that of every other coordinate, and the $O_v$ are
pairwise disjoint. The two inclusions are immediate: on $O_v$, $g_vg_x<0$
for $x\ne v$, and $g_xg_y>0$ for $x,y\ne v$.

(b) On $O_v$ every $x\ne v$ has the sign opposite to that of $v$, so the
option bad at $v$ is good at $x$.

(c) $O_v$ is an orthant of $\R^n$ contained in $H$. Conditioned on it, the
standard Gaussian vector has independent moduli $|g_1|,\dots,|g_n|$, each
with the half-normal density, which is positive on $(0,\infty)$. Put
$\xi_1\coloneqq|g_v|$, $\xi_2\coloneqq|g_j|$, $\xi_3\coloneqq|g_k|$. The map
$(\xi_1,\xi_2,\xi_3)\mapsto(\xi_1/\xi_2,\xi_1/\xi_3,\xi_1)$ is a
diffeomorphism of $(0,\infty)^3$ onto itself, with inverse
$(\rho_1,\rho_2,\xi_1)\mapsto(\xi_1,\xi_1/\rho_1,\xi_1/\rho_2)$, whose
Jacobian determinant has absolute value $\xi_1^2/(\rho_1^2\rho_2^2)>0$. So
$(t^v_j,t^v_k,|g_v|)$ has a positive density on $(0,\infty)^3$, and
integrating out the last coordinate shows that $(t^v_j,t^v_k)$ has a
positive density on $(0,\infty)^2$; the second marginal statement follows.
Nonempty open sets are non-null because $\mu$ has a positive density on
$H$.

(d) On $D_{uv}\cap\{t^u_v>\alpha\}$, $E$ gives credit at $u$; on
$O_v\subseteq D_{uv}$ the option good at $u$ is the option bad at $v$,
which is $\omega_v$. On $D_{uv}\cap\{t^u_v<\beta\}$, $E$ gives debt at $u$,
which on $O_u$ is $\omega_u$.

(e) By (a), the sets $Q_y$ are pairwise disjoint, and by (b), each
$\omega_y$ is credit at $x\notin\mathcal X$ on $Q_y$; $M$ has world-$x$
rank $1$ there. Off $\bigcup_yQ_y$ the two models agree.

(f) On $D_{xy}$ the numbers $e_x\cdot g=g_x$ and $e_y\cdot g=g_y$ have
opposite signs, so neither option has positive expectation under every
$p\in K$. If $x\in S(K)$, choose $y\in S(K)\setminus\{x\}$; then
$O_x\subseteq D_{xy}$ by (a).
\end{proof}

In Section~\ref{sec:LH}, sets defined by sign conditions and strict ratio
conditions on three coordinates are shown to be nonempty by exhibiting a
point: we prescribe the moduli of the three coordinates and their relative
signs, take the remaining coordinates nonzero and arbitrary, and choose the
overall sign so that $g\in H$; replacing $g$ by $-g$ changes no sign
condition of the form $g_ag_b<0$ and no ratio. Such sets are open, hence
non-null by (c).

\subsection{Results used}\label{sec:inputs}

We use the following results of FSE, with FSE's numbering. Each is stated
in FSE for every $n\ge3$: Lemmas~5, 6 and~11 and Proposition~13 are in
FSE's Sections~1.3 and~2, where the number of worlds is arbitrary, and
Lemma~8, Corollary~14 and Lemma~24 are stated for $n\ge3$.
\begin{enumerate}[label=(F\arabic*),leftmargin=*,itemsep=3pt]
\item\label{in:uf} (Lemma~5(a)) Every ultrafilter contains sets of
arbitrarily small measure. (Proof of Lemma~5(b)) If $\Iid$ is an ideal and
$R\notin\Iid$, some ultrafilter contains $R$ and avoids $\Iid$. The proof
in FSE is written for $\Jid_A$ but uses only that it is an ideal: the
family $\{R\}\cup\{H'\setminus J:J\in\Iid\}$ has the finite intersection
property, since $R\setminus(J_1\cup\dots\cup J_s)$ is non-null when
$R\notin\Iid$, and any ultrafilter extending it avoids $\Iid$ (the
three-world note, (F1)).
\item\label{in:ends} (Lemma~6) Assume (G) at worlds $v$ and $w$ and (S),
and write $C_q=(\alpha_q,\beta_q)$ in world $v$'s frame for
$q\in[e_v,e_w]$. (a) If $C_{q^*}=D_w$ and $q_m\to q^*$ along the edge,
then $\beta_{q_m}\to\infty$; if $C_{q^*}=D_v$ and $q_m\to q^*$, then
$\alpha_{q_m}\to0$. (b) $S_w\coloneqq\{q\in[e_v,e_w]:C_q=D_w\}$ is closed,
contains $e_w$ and not $e_v$; if $q_e$ is its point nearest to $e_v$,
every $q\in[e_v,q_e)$ has $C_q\ne D_w$. (c)
$S_v\coloneqq\{q\in[e_v,e_w]:C_q=D_v\}$ is closed, contains $e_v$ and not
$e_w$; it has a boundary point $q_0$ in the edge, and every boundary point
is a limit of points $q_m$ of the open edge with $C_{q_m}\ne D_v$.
\item\label{in:yard} (Lemma~8) Let $n\ge3$, assume (G) at every world,
(IB) at world~$v$ and (S), and let $w\ne v$. (a) If world $v$ is complete,
then for every non-null $A\subseteq H'$ there is $b<\infty$ such that
every measurable subset of $\Hi vwb$ is $\preceq^v_dA$; and for every open
interval $I\subseteq(0,\infty)$ there is a non-null $Z\subseteq D_{vw}\cap
\{t^v_w\in I\}$ with $Z\preceq^v_dA$. (b) For every non-null $T\subseteq
H'$ there is $a>0$ such that every measurable subset of $\Lo vwa$ is
$\preceq^v_cT$.

[In FSE's notation, with $v=1$ and $j\coloneqq w$: Lemma~8(a) gives
$\mathrm B_w(b,\infty)\prec_dA_0\preceq_dA$, subsets of the tail are
$\preceq_d$ the tail by weak (G), and $Z\coloneqq Z_I(A)$; Lemma~8(b) gives
$\mathrm B_w(0,a)\preceq_cT$. Lemma~8 holds for \emph{every} $j\ge2$, so
the edge $[e_v,e_w]$ may be prescribed.]
\item\label{in:SAC} (Lemma~11) Let world $v$ be complete with (G) and (IB),
and assume (G) at world~$w$ and (S). Let $C_{q^*}=D_w$ and let $q_m\to
q^*$ along $[e_v,e_w]$ with $\alpha_m\coloneqq\alpha^{(v)}_{q_m}<\infty$
for all $m$. Then for every $c<\infty$ and every non-null
$A'\subseteq D_{vw}\cap\{t^v_w<c\}$ there is $m_0$ such that for all $m\ge
m_0$ and every non-null $W\subseteq\Hi vw{\alpha_m}$ the pair $(A',W)$ is
strictly AC at~$v$.
\item\label{in:ext} (Proposition~13 and Corollary~14(a)) Let $n\ge3$,
assume (G) at every world and (S), and let world $v$ be complete with (IB).
If world $v$ is fine on $D_{vw}$, there are $\kappa_v\in(0,\infty)$ and
finitely additive probabilities $P^v_d,P^v_c$ on the measurable subsets of
$D_{vw}$ such that, with
$\Lambda_v(M)\coloneqq P^v_d(\debt_v(M)\cap D_{vw})-\kappa_vP^v_c(\cred_v(M)
\cap D_{vw})$:
(ii) $P^v_d(B)=0$ if and only if $B$ is null, and likewise for $P^v_c$;
(iii) for every $\eta>0$, $D_{vw}$ is a finite disjoint union of
measurable sets $Q$ with $P^v_d(Q)<\eta$ and $P^v_c(Q)<\eta$;
(iv) for regular $X,Y$ that agree off $D_{vw}$,
$\Lambda_v(X)<\Lambda_v(Y)$ implies $X\prec_vY$.
\item\label{in:EPfine} (Lemma~24) Let $n\ge3$, assume (G) at every world
and (S), and let worlds $v$ and $w$ be complete with (IB). If (EP$_{vw}$)
holds, the edge $[e_v,e_w]$ is fine.
\end{enumerate}
From CCF we use the following, with CCF's numbering.
\begin{enumerate}[label=(C\arabic*),leftmargin=*,itemsep=3pt]
\item\label{in:edge} (Lemmas~3 and~4) Assume (G) at worlds $v$ and $w$,
(S) and (EP$_{vw}$), and put $M_s\coloneqq E_w(s)$ for $s\in[0,\infty]$,
in world $v$'s frame. There is a continuous map $s\colon[e_v,e_w]\to
[-\infty,\infty]$ such that $C_q=M_{s(q)}$ whenever $s(q)\ge0$. If $s^*\in
(0,\infty)$ and $s_m\to s^*$ with $s_m<s^*$ for all $m$, there is $q^*$
with $s(q^*)=s^*$ such that every neighbourhood $N$ of $q^*$ has the
following property: for some $m_0$, every $s\in[s_{m_0},s^*]$ equals
$s(q)$ for some $q\in N$.
\item\label{in:reduction} (Theorem~16) Let $n\ge3$, let every ranking be
complete with (G) and (IB), and assume (S). If (EP$_{ij}$) holds for every
edge, then every coherent model whose credal set contains three vertices
is dominated.
\end{enumerate}
From the three-world note we use Definition~1 (connecting chains, their
perturbation sets, band-dense families) and the following two results,
with that note's numbering. Its Sections~3 and~4, which contain them, are
written for an arbitrary number $n\ge3$ of worlds.
\begin{enumerate}[label=(N\arabic*),leftmargin=*,itemsep=3pt]
\item\label{in:cut} (Theorem~3, no saturated cut) Let world $v$ be
complete, and assume (G) at worlds $v$ and $w$ and (S). For every
band-dense family $\Bfam$ of measurable subsets of $D_{vw}$ there is a
connecting chain from $D_v$ to $D_w$ with perturbation sets in $\Bfam$.
Every connecting chain from $D_v$ to $D_w$ has length at least $2$.
\item\label{in:count} (Lemma~4, counting lemma) Let $\preceq_v$ be a
preorder satisfying weak (G) and (IB) at world $v$. Let
$D_v=Y_0,X_1,Y_1,\dots,Y_{\ell-1},X_\ell=D_w$ be a connecting chain of
length $\ell\ge2$ whose perturbation sets $B_1,\dots,B_{\ell-1}$ are
contained in a measurable set $F$. Then there are no pairwise disjoint
measurable sets $P_1,\dots,P_\ell\subseteq D_{vw}\setminus F$ with
$D_v+\debt_v(B_i)\prec_vD_v+\debt_v(P_m)$ for all $1\le i\le\ell-1$ and
$1\le m\le\ell$.
\end{enumerate}

\emph{Standing hypotheses} (P1$_n$): $n\ge3$; every ranking is a complete
preorder satisfying (G) and (IB); (S). Each result below states which of
them it uses. Where a coherent model $W_0=D_K$ appears, it is assumed that
$|S(K)|\ge3$; by Lemma~\ref{lem:regions}(f), $W_0$ is silent on $O_x$ for
every $x\in S(K)$. Comparisons at a single world are stated with base $V$;
by (IB$'$) they transfer to any base model that is silent on the sets
involved, in particular to $W_0$ when these sets lie in regions $O_x$ with
$x\in S(K)$.

\section{Low-cone heaviness in \texorpdfstring{$n$}{n} worlds}\label{sec:LH}

In this section we prove FSE's Theorem~56 for every $n\ge3$. Nothing is
assumed about $V$ or about any other model.

\begin{definition}\label{def:CN}
(FSE, Definition~50.) At world $v$, a measurable $R\subseteq H'$ is
\emph{credit-null} if every pair $(A,T)$ with $A$ non-null and $T\subseteq
R$ is AC, and \emph{debt-infinitesimal} if every pair $(A,S)$ with
$A\subseteq R$ non-null and $S$ non-null is SI. $\CN_v$ and $\DI_v$ denote
the families of such sets; both are hereditary and invariant under null
modifications. We write $\CN_v\ne0$ if some non-null set is credit-null,
and likewise for $\DI_v$.
\end{definition}

\begin{lemma}\label{lem:excl}
(FSE, Lemma~51.) No world has both $\CN_v\ne0$ and $\DI_v\ne0$.
\end{lemma}
\begin{proof}
Let $L\in\CN_v$ and $R\in\DI_v$ be non-null. If $R\setminus L$ is
non-null, put $A\coloneqq R\setminus L$ and $T\coloneqq L$; otherwise
$R\subseteq L$ up to a null set, and we split $R$ into disjoint non-null
sets $A,T$. In both cases $(A,T)$ is a pair with $A\subseteq R$ and $T$
non-null, so it is SI because $R\in\DI_v$; and $A$ is non-null and
$T\subseteq L$ up to a null set, so it is AC because $L\in\CN_v$.
\end{proof}

\begin{lemma}\label{lem:conestails}
Let $n\ge3$, assume (G) at every world and (S), let world $v$ be complete
with (IB), and let $w\ne v$.
\begin{enumerate}[label=(\alph*),leftmargin=*,itemsep=2pt]
\item If $\CN_v\ne0$, there is $a>0$ with $\Lo vwa\in\CN_v$.
\item If $\DI_v\ne0$, there is $b<\infty$ with $\Hi vwb\in\DI_v$.
\end{enumerate}
\end{lemma}
\begin{proof}
(a) Let $Q\in\CN_v$ be non-null, and split it into disjoint non-null sets
$Q_1,Q_2$. By \ref{in:yard}(b) at world $v$ for the edge $[e_v,e_w]$ there
are $a_i>0$ such that every measurable subset of $\Lo vw{a_i}$ is
$\preceq^v_cQ_i$; put $a\coloneqq\min(a_1,a_2)$. Let $(A,T)$ be a pair with
$A$ non-null and $T\subseteq\Lo vwa$. Since $Q_1\cap Q_2=\varnothing$, $A$
is the union of $A\setminus Q_1$ and $A\setminus Q_2$, so one of them, say
$A'\coloneqq A\setminus Q_i$, is non-null. Then $(A',Q_i)$ is a pair with
non-null debt part and credit part in $Q\in\CN_v$, hence AC, and
$T\preceq^v_cQ_i$. The model $V+\debt_v(A')$ is silent on $Q_i\cup T$, since
$A'$ misses $Q_i$ and $A'\subseteq A$ misses $T$. So
\[
 V\preceq_vV+\debt_v(A')+\cred_v(Q_i)\preceq_vV+\debt_v(A')+\cred_v(T)
 \preceq_vV+\debt_v(A)+\cred_v(T).
\]
The first step says that $(A',Q_i)$ is AC (world $v$ is complete), the
second is $T\preceq^v_cQ_i$ in the context $V+\debt_v(A')$, by (IB$'$), and
the third is weak (G) (more debt). So $(A,T)$ is AC.

(b) Let $R\in\DI_v$ be non-null, and split it into disjoint non-null sets
$R_1,R_2$. By \ref{in:yard}(a) at world $v$ for $[e_v,e_w]$ there are
$b_i<\infty$ such that every measurable subset of $\Hi vw{b_i}$ is
$\preceq^v_dR_i$; put $b\coloneqq\max(b_1,b_2)$. Let $B\subseteq\Hi vwb$
and $S$ be non-null and disjoint. If $S\setminus R_1$ is non-null, put
$i\coloneqq1$ and $S''\coloneqq S\setminus R_1$; otherwise $S\subseteq R_1$
up to a null set, so $S$ misses $R_2$, and we put $i\coloneqq2$ and
$S''\coloneqq S$. Then $S''$ is non-null and misses $B\cup R_i$,
$B\preceq^v_dR_i$, and $(R_i,S'')$ is SI because $R\in\DI_v$. So
\[
 V+\debt_v(B)+\cred_v(S)\preceq_vV+\debt_v(B)+\cred_v(S'')\preceq_v
 V+\debt_v(R_i)+\cred_v(S'')\prec_vV,
\]
by weak (G) (less credit), by $B\preceq^v_dR_i$ in the context
$V+\cred_v(S'')$, which is silent on $B\cup R_i$, and because $(R_i,S'')$
is SI. So $(B,S)$ is SI.
\end{proof}

\begin{lemma}\label{lem:EPnotCN}
Let world $v$ be complete with (G) and (IB), assume (G) at world $w$ and
(S), and suppose (EP$_{vw}$). Then no low cone $\Lo vwa$, $a>0$, is
credit-null at~$v$.
\end{lemma}
\begin{proof}
This is FSE's Lemma~53, whose proof uses only the edge $[e_v,e_w]$. Fix
$0<b<a$ and a non-null $Z\subseteq D_{vw}\cap\{b<t^v_w<a\}$. Let $s$ and
$q^*$ be as in \ref{in:edge}, with $s^*\coloneqq b$ and
$s_m\coloneqq b-b/(m+1)$; so $C_{q^*}=E_w(b)$. The model $E_w(b)$ is
silent on $Z$, so $Y\coloneqq E_w(b)+\cred_v(Z)\prec_vE_w(b)$ by strict (G),
and by (T1), at the complete world $v$, $Y\prec_vC_q$ for all $q$ in some
neighbourhood $N$ of $q^*$. By \ref{in:edge} there are $s\in(0,b)$ and
$q\in N$ with $C_q=E_w(s)$; so $Y\prec_vE_w(s)$. Put
$L\coloneqq D_{vw}\cap\{s<t^v_w<b\}$, a non-null set disjoint from $Z$. Up to
the null level set $\{t^v_w=s\}$, $Y=E_w(s)+\debt_v(L)+\cred_v(Z)$, and
$E_w(s)$ is silent on $L\cup Z$. By (IB$'$),
$V+\debt_v(L)+\cred_v(Z)\prec_vV$: the pair $(L,Z)$ is SI, its debt part
is non-null and its credit part lies in $\Lo vwa$. So $\Lo vwa$ is not
credit-null.
\end{proof}

The next two lemmas are FSE's Lemmas~54 and~55. FSE's proofs use the
regions $P_v,P_w$ of three worlds; here they are replaced by the
three-world regions $D_{wv}\cap D_{wu}$ and $D_{vw}\cap D_{vu}$ for an
auxiliary world $u$. These are not regions of the form $O_x$ when $n\ge4$.
They may be used here because no dominance argument is involved: apart
from comparisons at a single world, the only step that concerns two
worlds is the competitiveness of a selection on its own two worlds.

\begin{lemma}\label{lem:CNtoDI}
Let $n\ge3$, assume (G) at every world and (S), and let worlds $v\ne w$ be
complete with (IB). Let $E$ be a selection on the open edge $(e_v,e_w)$
with $\beta^{(v)}_E>0$, and let $R\subseteq\cred_v(E)\cap D_{vw}$ be
credit-null at $v$. Then $R$ is debt-infinitesimal at $w$.
\end{lemma}
\begin{proof}
Fix a world $u\notin\{v,w\}$. Let $A\subseteq R$ and $S$ be non-null and
disjoint. Choose a non-null $S_0\subseteq S$ such that $S'\coloneqq
S\setminus S_0$ is non-null. By \ref{in:yard}(b) at world $w$ for the edge
$[e_w,e_u]$ there is $a'>0$ such that every measurable subset of $\Lo
wu{a'}$ is $\preceq^w_cS'$. Put
\[
 Q\coloneqq D_{vw}\cap D_{wu}\cap\{t^v_w<\beta^{(v)}_E\}\cap\{t^w_u<a'\}.
\]
$Q$ is open, and nonempty: take $|g_w|=1$, $|g_v|=\beta^{(v)}_E/2$ and
$|g_u|=2/a'$, with $g_v,g_u$ of one sign and $g_w$ of the other, the
remaining coordinates and the overall sign as in
Section~\ref{sec:regions}. So $Q$ is
non-null (Lemma~\ref{lem:regions}(c)). $E$ gives debt at $v$ on
$D_{vw}\cap\{t^v_w<\beta^{(v)}_E\}\supseteq Q$, so $Q$ misses
$A\subseteq\cred_v(E)$; and $Q\subseteq\Lo wu{a'}$, so $Q\preceq^w_cS'$.

Let $X$ be the model $E$ silenced on $Q\cup A$; so
$E=X+\debt_v(Q)+\cred_v(A)$. Split $Q$ into disjoint non-null sets
$Q_1,Q_2$. The pair $(Q_1,A)$ has non-null debt part and credit part
$A\subseteq R\in\CN_v$, so it is AC, and by strict (G) (more debt on $Q_2$)
\[
 V\preceq_vV+\debt_v(Q_1)+\cred_v(A)\prec_vV+\debt_v(Q)+\cred_v(A).
\]
By (IB$'$), $X\prec_vE$, so $E\preceq_vX$ fails; as $E$ is competitive on
$\{v,w\}$, $E\preceq_wX$.
By \eqref{eq:swap}, since $Q\cup A\subseteq D_{vw}$,
$E=X+\debt_w(A)+\cred_w(Q)$, and (IB$'$) gives
$V+\debt_w(A)+\cred_w(Q)\preceq_wV$. The model $V+\debt_w(A)$ is silent on
$Q\cup S'$, since $A$ misses $Q$ and $S$. So
\[
 V+\debt_w(A)+\cred_w(S)\prec_wV+\debt_w(A)+\cred_w(S')\preceq_w
 V+\debt_w(A)+\cred_w(Q)\preceq_wV,
\]
the first step by strict (G) (more credit on $S_0$, which misses $A$), the
second by $Q\preceq^w_cS'$ in the context $V+\debt_w(A)$. So $(A,S)$ is SI
at $w$.
\end{proof}

\begin{lemma}\label{lem:DItoCN}
Let $n\ge3$, assume (G) at every world and (S), and let worlds $v\ne w$ be
complete with (IB). If $\DI_w\ne0$, then $\CN_v\ne0$.
\end{lemma}
\begin{proof}
Fix a world $u\notin\{v,w\}$. By Lemma~\ref{lem:conestails}(b) at world $w$ for the edge $[e_w,e_v]$
there is $b<\infty$ with $\Hi wvb\in\DI_w$. On the edge $[e_v,e_w]$, in
world $v$'s frame, \ref{in:ends}(c) gives a boundary point $q_0$ of
$S_v=\{q:C_q=D_v\}$, so that $C_{q_0}=D_v$ as $S_v$ is closed, and points
$q_m\to q_0$ of the open edge with $C_m\coloneqq C_{q_m}\ne D_v$; by
\ref{in:ends}(a), $\alpha^{(v)}_m\coloneqq\alpha^{(v)}_{q_m}\to0$. In
world $w$'s frame, $C_m$ has thresholds
$\alpha^{(w)}_m=1/\beta^{(v)}_m$ and $\beta^{(w)}_m=1/\alpha^{(v)}_m$, which
tends to $\infty$.

\emph{Case I: $\beta^{(v)}_m>0$ for infinitely many $m$.} Along these $m$,
$\alpha^{(w)}_m<\infty$, and \ref{in:SAC} applies in world $w$'s frame on
$[e_w,e_v]$, at the point $q_0$, where $C_{q_0}=D_v$ is the far vertex
model. With $c\coloneqq b+1$ and the non-null set
$A'\coloneqq\Hi wvb\cap\{t^w_v<b+1\}$ it gives $m_0$ such that $(A',W)$ is
strictly AC at $w$ for every $m\ge m_0$ in this subsequence and every
non-null $W\subseteq\Hi wv{\alpha^{(w)}_m}$. Since
$\alpha^{(w)}_m\ge\beta^{(w)}_m\to\infty$, we may take such $m$ with
$\alpha^{(w)}_m>b+1$ and $W\coloneqq\Hi wv{\alpha^{(w)}_m}$, which is
non-null and disjoint from $A'$. But $A'\subseteq\Hi wvb\in\DI_w$, so
$(A',W)$ is SI: a contradiction. So Case~I does not occur.

\emph{Case II: $\beta^{(v)}_m=0$ for all large $m$.} Then
$C_m=(\alpha^{(v)}_m,0)$ with $0<\alpha^{(v)}_m<\infty$ for large $m$
($\alpha^{(v)}_m>0$ because $C_m\ne D_v$), and $C_m$ is silent on
$D_{vw}\cap\{t^v_w<\alpha^{(v)}_m\}=\Hi wv{1/\alpha^{(v)}_m}$. Fix such $m$
with $1/\alpha^{(v)}_m>b$, write $\alpha\coloneqq\alpha^{(v)}_m$, and put
\[
 T'\coloneqq D_{wv}\cap D_{wu}\cap\{t^v_w<\alpha\}.
\]
$T'$ is open and nonempty ($|g_w|=1$, $|g_v|=\alpha/2$, $|g_u|=1$, with
$g_v,g_u$ of one sign and $g_w$ of the other, the rest as in
Section~\ref{sec:regions}), hence non-null; it lies in
the silence of $C_m$; and $T'\subseteq\Hi wv{1/\alpha}\subseteq\Hi wvb$, so
$T'\in\DI_w$. We show $T'\in\CN_v$. A pair whose debt part is non-null and
whose credit part is null is AC by strict (G). So let $T\subseteq T'$ and
$A$ be non-null and disjoint. By \ref{in:yard}(a) at world $v$ for the edge
$[e_v,e_u]$ there is $b_A<\infty$ such that every measurable subset of $\Hi
vu{b_A}$ is $\preceq^v_dA$. Put
\[
 S\coloneqq D_{vw}\cap D_{vu}\cap\{t^v_u>b_A\}\cap
 \{\alpha/4<t^v_w<\alpha/2\}.
\]
$S$ is open and nonempty ($|g_v|=1$, $|g_w|=3/\alpha$, $|g_u|=1/(2b_A)$,
with $g_w,g_u$ of one sign and $g_v$ of the other, the rest as in
Section~\ref{sec:regions}), hence non-null. It lies
in the silence of $C_m$, and $S\preceq^v_dA$. It misses $T'$: on
$D_{vw}\cap D_{vu}$ the worlds $w$ and $u$ have the same sign, so $S\cap
D_{wu}=\varnothing$, while $T'\subseteq D_{wu}$. Let
$Y\coloneqq C_m+\debt_w(T)+\cred_w(S)$. The pair $(T,S)$ is SI at $w$,
since $T\subseteq T'\in\DI_w$ and $S$ is non-null; $C_m$ is silent on
$T\cup S$, so $Y\prec_wC_m$ by (IB$'$), and $C_m\preceq_wY$ fails. As
$C_m$ is competitive on $\{v,w\}$, $C_m\preceq_vY$. By \eqref{eq:swap},
$Y=C_m+\cred_v(T)+\debt_v(S)$, so (IB$'$) gives
\[
 V\preceq_vV+\debt_v(S)+\cred_v(T)\preceq_vV+\debt_v(A)+\cred_v(T),
\]
the second step by $S\preceq^v_dA$ in the context $V+\cred_v(T)$, which is
silent on $S\cup A$ ($S$ misses $T'\supseteq T$, and $A$ misses $T$). So
$(A,T)$ is AC, and $T'$ is a non-null member of $\CN_v$.
\end{proof}

\begin{theorem}\label{thm:LH}
Let $n\ge3$, let every ranking be a complete preorder satisfying (G) and
(IB), and assume (S). Then at every world $v$:
\begin{enumerate}[label=(\roman*),leftmargin=*,itemsep=2pt]
\item (LH$_v$): for every non-null $R\subseteq H'$ there is an SI pair
$(A,T)$ of world $v$ with $A$ non-null and $T\subseteq R$;
\item (DH$_v$): for every non-null $R\subseteq H'$ there is an AC pair
$(A,S)$ of world $v$ with $A\subseteq R$ and $A,S$ non-null.
\end{enumerate}
No assumption on $V$ is used.
\end{theorem}
\begin{proof}
(i) says $\CN_v=0$ and (ii) says $\DI_v=0$. Suppose $\CN_v\ne0$ for some
$v$, and choose distinct worlds $w,u\ne v$. By
Lemma~\ref{lem:conestails}(a) for the edges $[e_v,e_w]$ and $[e_v,e_u]$,
taking the smaller of the two levels, there is $a>0$ with $\Lo vwa,\Lo
vua\in\CN_v$. By Lemma~\ref{lem:EPnotCN}, (EP$_{vw}$) and (EP$_{vu}$) fail:
there are selections $E$ on $[e_v,e_w]$ and $E'$ on $[e_v,e_u]$ whose
credal sets omit both endpoints of their edges; they lie in the open edges
and are two-sided. The sets
\[
 R\coloneqq D_{vw}\cap D_{vu}\cap\{t^v_w>\alpha^{(v)}_E\}\cap\{t^v_u<a\},
 \qquad
 R'\coloneqq D_{vw}\cap D_{vu}\cap\{t^v_u>\alpha^{(v)}_{E'}\}\cap
 \{t^v_w<a\}
\]
are open and nonempty (for $R$: $|g_v|=1$, $|g_w|=1/(2\alpha^{(v)}_E)$,
$|g_u|=2/a$, with $g_w,g_u$ of one sign and $g_v$ of the other, the rest as
in Section~\ref{sec:regions}; likewise for $R'$), hence non-null. $R\subseteq\Lo vua$ and $R'\subseteq\Lo vwa$, so
both are credit-null at $v$; and $R\subseteq\cred_v(E)\cap D_{vw}$,
$R'\subseteq\cred_v(E')\cap D_{vu}$. Since $\beta^{(v)}_E$ and
$\beta^{(v)}_{E'}$ are positive, Lemma~\ref{lem:CNtoDI} gives $R\in\DI_w$
and $R'\in\DI_u$. Lemma~\ref{lem:DItoCN}, applied to the worlds $u$ and $w$
in the roles of $v$ and $w$ there, turns $\DI_w\ne0$ into $\CN_u\ne0$.
Together with $\DI_u\ne0$ this contradicts Lemma~\ref{lem:excl} at
world~$u$. So $\CN_v=0$ at every world. If $\DI_w\ne0$ for some $w$,
Lemma~\ref{lem:DItoCN} gives $\CN_v\ne0$ for any $v\ne w$, which has just
been excluded.
\end{proof}

\begin{corollary}\label{cor:accept}
Under the hypotheses of Theorem~\ref{thm:LH}, let $w\ne v$. For every
non-null $B\subseteq H'$ there is $c<\infty$ such that $(Y,B)$ is an SI
pair of world $w$ for every non-null $Y\subseteq\Hi wvc$ disjoint from $B$.
In particular, if $u,v,w$ are distinct and $B\subseteq O_u$ is non-null,
then
\[
 V+\omega_w(Y)+\omega_u(B)\prec_wV\qquad\text{for every non-null }
 Y\subseteq O_w\cap\Hi wvc.
\]
\end{corollary}
\begin{proof}
By \ref{in:yard}(b) at world $w$ for the edge $[e_w,e_v]$ there is $a'>0$
such that every measurable subset of $\Lo wv{a'}$ is $\preceq^w_cB$. By
Theorem~\ref{thm:LH}(i) at $w$, with $R\coloneqq\Lo wv{a'}$, there is an SI
pair $(A,T'')$ of world $w$ with $A$ non-null and $T''\subseteq\Lo
wv{a'}$. By \ref{in:yard}(a) at world $w$ for $[e_w,e_v]$ there is $b$ such
that every measurable subset of $\Hi wvb$ is $\preceq^w_dA$. Put
$c\coloneqq\max(a',b)$, and let $Y\subseteq\Hi wvc$ be non-null and
disjoint from $B$. Then $Y$ misses $T''\subseteq\{t^w_v<a'\}$, and
\begin{multline*}
 V+\debt_w(Y)+\cred_w(B)\preceq_wV+\debt_w(Y)+\cred_w(T'')\\
 \preceq_wV+\debt_w(A)+\cred_w(T'')\prec_wV.
\end{multline*}
The first step is $T''\preceq^w_cB$ in the context $V+\debt_w(Y)$, which is
silent on $B\cup T''$; the second is $Y\preceq^w_dA$ in the context
$V+\cred_w(T'')$, which is silent on $Y\cup A$ ($Y$ and $A$ miss $T''$);
the third says that $(A,T'')$ is SI. For the last statement, $Y\subseteq
O_w$ and $B\subseteq O_u$ are disjoint, $\omega_w$ is debt at $w$ on $O_w$,
and $\omega_u$ is credit at $w$ on $O_u$ (Lemma~\ref{lem:regions}).
\end{proof}

This is FSE's Corollary~58(c) for every $n$. It says that, against credit
on any non-null set, world $w$ accepts debt on a deep enough tail of each
of its edges. In Sections~\ref{sec:own} and~\ref{sec:dichotomy} it lets the
third world of a triple accept the change that the other two need.

\section{The own-region theorem}\label{sec:own}

In this section and the next, (P1$_n$) holds, $K$ is the credal set of a
coherent model $W_0\coloneqq D_K$ with $|S(K)|\ge3$, and $W_0$ is silent on
$O_x$ for every $x\in S(K)$ (Lemma~\ref{lem:regions}(f)). For $K=\Delta$,
$W_0=V$ and $S(K)$ is the set of all worlds.

For a selection $E$ on an edge $[e_u,e_k]$, with thresholds $(\alpha,\beta)$
in world $u$'s frame, and measurable sets $T\subseteq O_k\cap
\{t^u_k>\alpha\}$ and $B\subseteq O_u\cap\{t^u_k<\beta\}$, put
\[
 M+\tau_{T,B}\coloneqq M+\omega_k(T)+\omega_u(B)
\]
for regular $M$. By Lemma~\ref{lem:regions}(d), $E$ gives $\omega_k$ on
$T$ and $\omega_u$ on $B$. At world $u$, $\tau_{T,B}$ is credit on $T$ and
debt on $B$; at world $k$ it is debt on $T$ and credit on $B$; at every
other world it is credit on $T\cup B$.

\begin{lemma}\label{lem:XOR}
Assume (IB) at worlds $u$ and $k$. Let $E$ be a selection on the open edge
$(e_u,e_k)$ with thresholds $(\alpha,\beta)$ in world $u$'s frame, let
$T\subseteq O_k\cap\{t^u_k>\alpha\}$ and $B\subseteq O_u\cap\{t^u_k<\beta\}$
be measurable, and let $M$ be a regular model silent on $T\cup B$. Then
$M+\tau_{T,B}\preceq_uM$ or $M+\tau_{T,B}\preceq_kM$.
\end{lemma}
\begin{proof}
Let $E'$ be the model $E$ silenced on $T\cup B$. The pairs $(E,E')$ and
$(M+\tau_{T,B},M)$ have the same difference set $T\cup B$, and the same
values on it: $\omega_k$ on $T$ and $\omega_u$ on $B$ in the first model,
silence in the second. By (IB$'$), $E\preceq_xE'$ if and only if
$M+\tau_{T,B}\preceq_xM$, for $x=u,k$. As $E$ is competitive on $\{u,k\}$,
$E\preceq_uE'$ or $E\preceq_kE'$.
\end{proof}

This is FSE's Lemma~25 for every $n$; we need only its first half. The
\emph{dominance steps} below add to $W_0$ the options $\omega_x$ on subsets
of regions $O_x$ with $x$ in a triple of worlds of $S(K)$; by
Lemma~\ref{lem:regions}(e) they then improve every world outside the triple
strictly.

\begin{theorem}\label{thm:own}
Assume (P1$_n$) and suppose that $W_0=D_K$ is undominated. Let $u,k,l$ be
distinct worlds in $S(K)$, and let $E$ be a two-sided selection on
$(e_u,e_k)$, with thresholds $(\alpha,\beta)$ in world $u$'s frame. Put
\[
 T_0\coloneqq O_k\cap\{t^u_k>\alpha\},\qquad
 B_0\coloneqq O_u\cap\{t^u_k<\beta\},\qquad
 Y_c\coloneqq O_l\cap\Hi lkc=O_l\cap\{t^k_l<1/c\}\quad(c>0).
\]
Then for every non-null $T\subseteq T_0$ there is $c_0\in(0,\infty)$ such
that, for every $c\ge c_0$, $T$ is not a finite union of measurable sets $Q$
for which $(Q,Y_c)$ is a weakly SI pair of world~$k$.
\end{theorem}

On $O_k$, $\omega_k$ is debt at $k$, and on $O_l$, $\omega_l$ is credit at
$k$. So the theorem says that world $k$'s debt on small pieces of $T$ is
not compensated by its credit on deep tails of $O_l$: world $k$'s debt
order is non-fine on its own region, relative to credit on deep tails of
the region of a third world. The sets $T_0$, $B_0$ and $Y_c$ are non-null
by Lemma~\ref{lem:regions}(c), and $E$ gives $\omega_k$ on $T_0$ and
$\omega_u$ on $B_0$. For $n=3$ and $(u,k,l)=(1,2,3)$ this is FSE's
Theorem~64.

\begin{proof}
If $c\ge c_0$, then $Y_c\subseteq Y_{c_0}$, and a weakly SI pair $(Q,Y_c)$
of world $k$ gives a weakly SI pair $(Q,Y_{c_0})$, by weak (G) (more
credit). So it suffices to find one $c_0$ for which the conclusion holds.
Suppose there is none:
\begin{enumerate}[label=(M2$^*$),leftmargin=*]
\item for every $c\in(0,\infty)$, $T$ is a finite union of measurable sets
$Q$ for which $(Q,Y_c)$ is a weakly SI pair of world~$k$.
\end{enumerate}

\emph{Step 1: for every non-null $Z\subseteq B_0$,
$W_0\prec_uW_0+\tau_{T,Z}$.} Let $Z\subseteq B_0$ be non-null. By
Corollary~\ref{cor:accept} at world $l$, for the edge $[e_l,e_k]$ and the
set $Z\subseteq O_u$, there is $c$ such that
$V+\omega_l(Y)+\omega_u(Z)\prec_lV$ for every non-null $Y\subseteq O_l\cap\Hi
lkc$. Put $Y\coloneqq Y_c$; it is non-null and disjoint from $Z$. By
(M2$^*$), $T=Q_1\cup\dots\cup Q_r$ with $(Q_i,Y)$ weakly SI at $k$;
replacing $Q_i$ by $Q_i\setminus(Q_1\cup\dots\cup Q_{i-1})$ keeps this, by
weak (G) (less debt), so we may take the $Q_i$ pairwise disjoint. Put
$\Sigma_i\coloneqq Q_1\cup\dots\cup Q_i$ for $0\le i\le r$, and suppose
that $W_0+\tau_{T,Z}\preceq_uW_0$. Since $\omega_u$ is debt at $u$ on the
non-null set $Z$, $W_0\prec_uW_0+\tau_{\Sigma_0,Z}$ by strict (G). So there
is a least $i^*\ge1$ with $W_0+\tau_{\Sigma_{i^*},Z}\preceq_uW_0$, and
$W_0+\tau_{\Sigma_{i^*-1},Z}\preceq_uW_0$ fails. By Lemma~\ref{lem:XOR},
with $M\coloneqq W_0$, which is silent on $O_k\cup O_u$,
$W_0+\tau_{\Sigma_{i^*-1},Z}\preceq_kW_0$. Let
\[
 d\coloneqq\omega_k(\Sigma_{i^*})+\omega_u(Z)+\omega_l(Y),
\]
a change on the pairwise disjoint sets $\Sigma_{i^*}\subseteq O_k$,
$Z\subseteq O_u$ and $Y\subseteq O_l$, on which $W_0$ is silent.
\begin{itemize}[leftmargin=*,itemsep=2pt]
\item \emph{World $u$.} $W_0+d$ is $W_0+\tau_{\Sigma_{i^*},Z}$ with credit
at $u$ added on $Y$. By strict (G),
$W_0+d\prec_uW_0+\tau_{\Sigma_{i^*},Z}\preceq_uW_0$.
\item \emph{World $k$.} The models $W_0+d$ and
$W_0+\tau_{\Sigma_{i^*-1},Z}$ differ exactly on $Q_{i^*}\cup Y$; their
world-$k$ ranks are $2$ and $1$ on $Q_{i^*}$, and $0$ and $1$ on $Y$. So do
$V+\debt_k(Q_{i^*})+\cred_k(Y)$ and $V$, and the first is $\preceq_k$ the
second because $(Q_{i^*},Y)$ is weakly SI. By (IB$'$),
$W_0+d\preceq_kW_0+\tau_{\Sigma_{i^*-1},Z}\preceq_kW_0$.
\item \emph{World $l$.} $\omega_k$ is credit at $l$ on $O_k$, so by weak
(G), $W_0+d\preceq_lW_0+\omega_l(Y)+\omega_u(Z)$; and
$W_0+\omega_l(Y)+\omega_u(Z)\prec_lW_0$ by the choice of $c$ and (IB$'$).
\item \emph{Every world $x\notin\{u,k,l\}$.} $W_0+d\prec_xW_0$ by
Lemma~\ref{lem:regions}(e).
\end{itemize}
So $W_0+d$ dominates $W_0$, a contradiction. Hence
$W_0+\tau_{T,Z}\preceq_uW_0$ fails, and completeness gives Step~1.

\emph{Step 2: a contradiction.} We show that $T$ is credit-null at
world~$u$, which Theorem~\ref{thm:LH} excludes. Let $A$ be non-null and
$T''\subseteq T$ disjoint from $A$, and suppose that $(A,T'')$ is SI at $u$.
By \ref{in:yard}(a) at world $u$ for the edge $[e_u,e_l]$ there is $b$ such
that every measurable subset of $\Hi ulb$ is $\preceq^u_dA$. The set
\[
 Z\coloneqq O_u\cap\{t^u_k<\beta\}\cap\{t^u_l>b\}\subseteq B_0
\]
is non-null (Lemma~\ref{lem:regions}(c)), and $Z\preceq^u_dA$. The model
$V+\cred_u(T'')$ is silent on $Z\cup A$, since $Z\subseteq O_u$ misses
$T''\subseteq O_k$ and $A$ misses $T''$. So
\[
 V+\debt_u(Z)+\cred_u(T)\preceq_uV+\debt_u(Z)+\cred_u(T'')\preceq_u
 V+\debt_u(A)+\cred_u(T'')\prec_uV,
\]
by weak (G) (more credit on $T\setminus T''$, which misses $Z$), by
$Z\preceq^u_dA$ in the context $V+\cred_u(T'')$, and because $(A,T'')$ is SI.
On $O_u$, $\omega_u$ is debt at $u$, and on $O_k$, $\omega_k$ is credit at
$u$; so this says $V+\tau_{T,Z}\prec_uV$, and by (IB$'$)
$W_0+\tau_{T,Z}\prec_uW_0$, contrary to Step~1. Hence every pair $(A,T'')$
with $A$ non-null and $T''\subseteq T$ is AC at $u$, and $T$ is a non-null
credit-null set at world~$u$.
\end{proof}

\section{The endpoint dichotomy}\label{sec:dichotomy}

\begin{proposition}\label{prop:partner}
Assume (P1$_n$) and suppose that $W_0=D_K$ is undominated. Let $u,v,w$ be
distinct worlds in $S(K)$, and suppose that world $v$ is fine on $D_{vw}$.
Then no selection on the open edge $(e_u,e_v)$ is two-sided. Hence
(EP$_{uv}$) holds, and the edge $[e_u,e_v]$ is fine.
\end{proposition}
\begin{proof}
By \ref{in:ext}, world $v$ has a local extension $\Lambda_v$ on $D_{vw}$,
with $P^v_d$, $P^v_c$ and $\kappa_v$. Suppose that $E$ is a two-sided
selection on $(e_u,e_v)$, with thresholds $(\alpha,\beta)$ in world $u$'s
frame, and put
\[
 T_0\coloneqq O_v\cap\{t^u_v>\alpha\},\qquad
 B_0\coloneqq O_u\cap\{t^u_v<\beta\},\qquad
 \tau_{T,B}\coloneqq\omega_v(T)+\omega_u(B)
\]
for $T\subseteq T_0$ and $B\subseteq B_0$. Both sets are non-null
(Lemma~\ref{lem:regions}(c); $0<\beta\le\alpha<\infty$), and $E$ gives
$\omega_v$ on $T_0$ and $\omega_u$ on $B_0$ (Lemma~\ref{lem:regions}(d)).

\emph{Step 1: for every non-null $B\subseteq B_0$ and every $T\subseteq T_0$,
$W_0\prec_uW_0+\tau_{T,B}$.} By Corollary~\ref{cor:accept} at world $w$,
for the edge $[e_w,e_v]$ and the set $B\subseteq O_u$, there is $c$ such
that $Z\coloneqq O_w\cap\Hi wvc$, which is non-null by
Lemma~\ref{lem:regions}(c), satisfies
$V+\omega_w(Z)+\omega_u(B)\prec_wV$. As $Z\subseteq O_w\subseteq D_{vw}$ is
non-null, $\eta\coloneqq\kappa_vP^v_c(Z)>0$ by \ref{in:ext}(ii). Suppose
that $W_0+\tau_{T^*,B}\preceq_uW_0$ for some $T^*\subseteq T_0$. By
\ref{in:ext}(iii), intersected with $T^*\subseteq O_v\subseteq D_{vw}$,
$T^*=Q_1\cup\dots\cup Q_r$ with pairwise disjoint measurable $Q_i$ and
$P^v_d(Q_i)<\eta$. Put $\Sigma_i\coloneqq Q_1\cup\dots\cup Q_i$. As
$W_0\prec_uW_0+\tau_{\Sigma_0,B}$ by strict (G) (debt at $u$ on $B$), there
is a least $i^*\ge1$ with $W_0+\tau_{\Sigma_{i^*},B}\preceq_uW_0$, and
$W_0+\tau_{\Sigma_{i^*-1},B}\preceq_uW_0$ fails. By Lemma~\ref{lem:XOR},
applied to the edge $(e_u,e_v)$ with $M\coloneqq W_0$, which is silent on
$O_v\cup O_u$, $W_0+\tau_{\Sigma_{i^*-1},B}\preceq_vW_0$. Let
\[
 d\coloneqq\omega_v(\Sigma_{i^*})+\omega_u(B)+\omega_w(Z),
\]
a change on the pairwise disjoint sets $\Sigma_{i^*}\subseteq O_v$,
$B\subseteq O_u$ and $Z\subseteq O_w$, on which $W_0$ is silent.
\begin{itemize}[leftmargin=*,itemsep=2pt]
\item \emph{World $u$.} $W_0+d\prec_uW_0+\tau_{\Sigma_{i^*},B}\preceq_uW_0$
by strict (G): $\omega_w$ is credit at $u$ on $Z$.
\item \emph{World $v$.} The change $\delta\coloneqq\omega_v(Q_{i^*})+
\omega_w(Z)$ is debt at $v$ on $Q_{i^*}$ and credit at $v$ on $Z$, and both
sets lie in $D_{vw}$. The models $V+\delta$ and $V$ agree off $D_{vw}$,
$\Lambda_v(V)=0$ and
$\Lambda_v(V+\delta)=P^v_d(Q_{i^*})-\kappa_vP^v_c(Z)<0$; so
$V+\delta\prec_vV$ by \ref{in:ext}(iv). The pairs
$(W_0+d,W_0+\tau_{\Sigma_{i^*-1},B})$ and $(V+\delta,V)$ differ exactly on
$Q_{i^*}\cup Z$, with the same values there; by (IB$'$),
$W_0+d\prec_vW_0+\tau_{\Sigma_{i^*-1},B}\preceq_vW_0$.
\item \emph{World $w$.} $\omega_v$ is credit at $w$ on $O_v$, so
$W_0+d\preceq_wW_0+\omega_w(Z)+\omega_u(B)\prec_wW_0$, by weak (G) and by
the choice of $Z$ and (IB$'$).
\item \emph{Every world $x\notin\{u,v,w\}$.} $W_0+d\prec_xW_0$ by
Lemma~\ref{lem:regions}(e).
\end{itemize}
So $W_0+d$ dominates $W_0$, a contradiction. Hence
$W_0+\tau_{T,B}\preceq_uW_0$ fails for every $T\subseteq T_0$, and
completeness gives Step~1.

\emph{Step 2: $T_0$ is credit-null at world $u$.} Let $(A,T)$ be a pair of
world $u$ with $A$ non-null and $T\subseteq T_0$, and suppose that it is
SI. By \ref{in:yard}(a) at world $u$ for the edge $[e_u,e_w]$ there is $b$
such that every measurable subset of $\Hi uwb$ is $\preceq^u_dA$. The set
$Z'\coloneqq O_u\cap\{t^u_v<\beta\}\cap\{t^u_w>b\}\subseteq B_0$ is
non-null (Lemma~\ref{lem:regions}(c)), and $Z'\preceq^u_dA$. The model $V+\cred_u(T)$ is silent on
$Z'\cup A$ ($Z'\subseteq O_u$ misses $T\subseteq O_v$, and $A$ misses $T$),
so
\[
 V+\debt_u(Z')+\cred_u(T)\preceq_uV+\debt_u(A)+\cred_u(T)\prec_uV,
\]
that is, $V+\tau_{T,Z'}\prec_uV$, and by (IB$'$)
$W_0+\tau_{T,Z'}\prec_uW_0$, contrary to Step~1. So $T_0$ is a non-null
credit-null set at world $u$, which Theorem~\ref{thm:LH} excludes.

Hence no selection on $(e_u,e_v)$ is two-sided. A selection on
$[e_u,e_v]$ whose credal set omits both endpoints lies in the open edge
and is two-sided; so (EP$_{uv}$) holds, and \ref{in:EPfine} makes the edge
$[e_u,e_v]$ fine.
\end{proof}

This is FSE's Lemmas~26--29 in $n$ worlds, with $P_1,P_2,P_3$ replaced by
$O_u,O_v,O_w$. FSE's Lemma~26 uses the fineness of world $w$ on $D_{vw}$;
here Corollary~\ref{cor:accept}, that is, low-cone heaviness at $w$, takes
its place, so only world $v$ needs to be fine on $D_{vw}$. This is what
lets the endpoint property spread along edges one at a time.

\begin{theorem}\label{thm:dichotomy}
Assume (P1$_n$) and suppose that $W_0=D_K$ is undominated. Then either
(EP$_{xy}$) holds for all distinct $x,y\in S(K)$, or it holds for no such
pair.
\end{theorem}
\begin{proof}
Suppose (EP$_{vw}$) holds for some distinct $v,w\in S(K)$. By
\ref{in:EPfine}, worlds $v$ and $w$ are fine on $D_{vw}$. Let $u\in
S(K)\setminus\{v,w\}$. Proposition~\ref{prop:partner} with $(u,v,w)$ gives
(EP$_{uv}$), and with $(u,w,v)$, world $w$ being fine on $D_{wv}$, it gives
(EP$_{uw}$). Let $x,y\in S(K)\setminus\{v,w\}$ be distinct. By what was just
shown (EP$_{xv}$) holds, so world $x$ is fine on $D_{xv}$ by
\ref{in:EPfine}, and Proposition~\ref{prop:partner} with $(u,v,w)\coloneqq
(y,x,v)$ gives (EP$_{yx}$).
\end{proof}

\section{Rich far ends}\label{sec:rich}

In this section and the next, $1$, $2$ and $3$ denote three distinct worlds;
by relabelling they may be any three distinct worlds (FSE, Section~1.3).
We work in world~$1$'s frame on the edge $[e_1,e_2]$:
$t\coloneqq t^1_2$ on $D_{12}$, and $C_q=(\alpha_q,\beta_q)\coloneqq
(\alpha^{(1)}_q,\beta^{(1)}_q)$ for $q\in[e_1,e_2]$. Let $q_e$ be the point
of $S_2=\{q\in[e_1,e_2]:C_q=D_2\}$ nearest to $e_1$ (\ref{in:ends}(b));
$q_e\ne e_1$. As in the three-world note, Section~5, $q_e$ is the \emph{far
end of the frame $(1,2)$}. We parametrize the edge by
$q(s)\coloneqq(1-s)e_1+se_2$, $s\in[0,1]$, write $C(s)\coloneqq C_{q(s)}$,
with thresholds $\alpha(s),\beta(s)$, and let $s_e\in(0,1]$ be the
parameter of $q_e$. For $0<s<s_e$, $q(s)$ lies in the open edge and
$C(s)\ne D_2$. In world~$1$'s frame, $E_2(b)=(\infty,b)$ gives debt at $1$
on $D_{12}\cap\{t<b\}$, silence on the rest of $D_{12}$, and the option
good at $1$ and $2$ off $D_{12}$.

\begin{definition}\label{def:rich}
(The three-world note, Definition~5.) The far end of the frame $(1,2)$ is
\emph{rich} if there are $s_m\in(0,s_e)$ with $s_m\to s_e$ and
$\alpha(s_m)<\infty$ for every $m$.
\end{definition}

In world~$2$'s frame, $\alpha(s)<\infty$ means
$\beta^{(2)}_{q(s)}=1/\alpha(s)>0$: the selection $C(s)$ gives debt at
world~$2$ on the non-null set $D_{21}\cap\{t^2_1<1/\alpha(s)\}$.

\begin{proposition}\label{prop:rich}
Assume (P1$_n$), and suppose that $W_0=D_K$ is undominated, that
$1,2,3\in S(K)$, and that the edge $[e_1,e_3]$ carries a two-sided
selection. Then the far end of the frame $(1,2)$ is rich.
\end{proposition}
\begin{proof}
We follow the proof of the three-world note's Proposition~6. Its
Steps~1--4 use only the edge $[e_1,e_2]$, (G) at every world, completeness
and (IB) at world~$1$, and (S); they are repeated here because that note
states them for $n=3$. Its Step~5 is replaced by an argument on the regions
$O_1$ and $O_2$.

Suppose the far end is not rich. Then there is $s_0\in(0,s_e)$ with
$\alpha(s)=\infty$ for all $s\in(s_0,s_e)$. For such $s$, $C(s)\ne D_2$, so
$\beta(s)<\infty$, and $C(s)=E_2(\beta(s))$. For $0\le b<b'<\infty$, the
models $E_2(b)$ and $E_2(b')$ differ exactly on the non-null band
$D_{12}\cap\{b\le t<b'\}$, where $E_2(b)$ is silent and $E_2(b')$ gives
debt at $1$, that is, credit at $2$ by \eqref{eq:swap}. By strict (G) at
worlds $1$ and $2$,
\begin{equation}
 E_2(b)\prec_1E_2(b')\quad\text{and}\quad E_2(b')\prec_2E_2(b)
 \qquad(0\le b<b'<\infty).
 \label{eq:E2}
\end{equation}

\emph{Step 1: $\beta$ is continuous on $(s_0,s_e)$, and $\beta(s)\to\infty$
as $s\uparrow s_e$.} Let $b\in[0,\infty)$ and $s\in(s_0,s_e)$. By
\eqref{eq:E2} and reflexivity, $C(s)\preceq_1E_2(b)$ holds if and only if
$\beta(s)\le b$, and $C(s)\preceq_2E_2(b)$ holds if and only if
$\beta(s)\ge b$. The sections $A_1(E_2(b))$ and $A_2(E_2(b))$ are closed by
(S). So $\{s\in(s_0,s_e):\beta(s)\le b\}$ and
$\{s\in(s_0,s_e):\beta(s)\ge b\}$ are closed in $(s_0,s_e)$ for every
$b\in[0,\infty)$, and $\beta$ is lower and upper semicontinuous on
$(s_0,s_e)$, hence continuous. Since $C_{q_e}=D_2$, \ref{in:ends}(a) gives
$\beta(s)\to\infty$ as $s\uparrow s_e$.

\emph{Step 2: every level above some $z_0$ is crossed from both sides.} Fix
$s_1\in(s_0,s_e)$ and put $z_0\coloneqq\beta(s_1)$. Let $z>z_0$. By Step~1
there is $s_2\in(s_1,s_e)$ with $\beta>z$ on $[s_2,s_e)$. The set
$\{s\in[s_1,s_2]:\beta(s)=z\}$ is nonempty by the intermediate value
theorem and compact by continuity; let $s^-$ and $s^+$ be its least and its
greatest element. Then $s_1<s^-\le s^+<s_2$, and, again by the intermediate
value theorem, $\beta<z$ on $[s_1,s^-)$ and $\beta>z$ on $(s^+,s_e)$.
Moreover $C(s^-)=C(s^+)=E_2(z)$.

\emph{Step 3: bands at levels above $z_0$.} Let $z>z_0$ and let
$A\subseteq H'$ be non-null. We show that there is $\varepsilon>0$ with
\begin{equation}
 D_{12}\cap\{z-\varepsilon<t<z\}\preceq_dA\quad\text{and}\quad
 D_{12}\cap\{z<t<z+\varepsilon\}\preceq_dA.
 \label{eq:bands}
\end{equation}

(i) By \ref{in:yard}(a) at world~$1$ for the edge $[e_1,e_2]$ there is a
non-null $Z\subseteq D_{12}\cap\{z/4<t<z/2\}$ with $Z\preceq_dA$. Let
$X\coloneqq E_2(z)+\sil(Z)$. As $E_2(z)$ gives debt at $1$ on $Z$,
$X\prec_1E_2(z)=C(s^-)$ by strict (G), and (T1), at the complete
world~$1$, gives $X\prec_1C_q$ for all $q$ near $q(s^-)$. By continuity
$\beta(s)\to z$ as $s\uparrow s^-$, so there is $s'\in[s_1,s^-)$ with
$X\prec_1C(s')$ and $b'\coloneqq\beta(s')$ in $(z/2,z)$. Put
$L\coloneqq D_{12}\cap\{b'<t<z\}$ and $X_0\coloneqq E_2(b')+\sil(Z)$. Since
$Z\subseteq\{t<z/2\}\subseteq\{t<b'\}$, we have
$C(s')=E_2(b')=X_0+\debt_1(Z)$, and, up to the null level set $\{t=b'\}$,
$X=X_0+\debt_1(L)$: both give debt at $1$ exactly on
$D_{12}\cap\{t<z\}\setminus Z$, silence on the rest of $D_{12}$, and the
option good at $1$ and $2$ off $D_{12}$. So
$X_0+\debt_1(L)\prec_1X_0+\debt_1(Z)$, where $X_0$ is silent on the disjoint
sets $L$ and $Z$. By (IB$'$), $V+\debt_1(L)\prec_1V+\debt_1(Z)$, that is,
$L\prec_dZ$; so $L\preceq_dZ\preceq_dA$.

(ii) By \ref{in:yard}(a) there is a non-null
$Z'\subseteq D_{12}\cap\{z+1<t<z+2\}$ with $Z'\preceq_dA$; $E_2(z)$ is
silent on $Z'$. Let $X'\coloneqq E_2(z)+\debt_1(Z')$. By \eqref{eq:swap},
$X'$ gives credit at $2$ on $Z'$, so $X'\prec_2E_2(z)=C(s^+)$ by strict (G)
at world~$2$. By (T2) with $(v,w)=(1,2)$, $C_q\preceq_1X'$ for all $q$ in
the open edge near $q(s^+)$. By continuity $\beta(s)\to z$ as $s\downarrow
s^+$, so there is $s''\in(s^+,s_e)$ with $C(s'')\preceq_1X'$ and
$b''\coloneqq\beta(s'')$ in $(z,z+1)$. Put $L'\coloneqq D_{12}\cap
\{z<t<b''\}$. Up to a null level set, $C(s'')=E_2(b'')=E_2(z)+\debt_1(L')$,
and $E_2(z)$ is silent on the disjoint sets $L'$ and $Z'$. So
$E_2(z)+\debt_1(L')\preceq_1E_2(z)+\debt_1(Z')$, and (IB$'$) gives
$L'\preceq_dZ'\preceq_dA$.

With $\varepsilon\coloneqq\min(z-b',b''-z)$, \eqref{eq:bands} follows from
(i), (ii) and weak (G).

\emph{Step 4: a band region in every $\Jid_A$.} Put
$R\coloneqq D_{12}\cap\{z_0+1<t<z_0+2\}$, and let $A\subseteq H'$ be
non-null. For each $z\in[z_0+1,z_0+2]$, Step~3 gives $\varepsilon_z>0$ such
that $D_{12}\cap\{|t-z|<\varepsilon_z\}$ is, up to the null set $\{t=z\}$,
the union of two sets $\preceq_dA$; so it belongs to $\Jid_A$. Finitely many
of the intervals $(z-\varepsilon_z,z+\varepsilon_z)$ cover the compact
interval $[z_0+1,z_0+2]$, and $\Jid_A$ is an ideal; so $R\in\Jid_A$. This
holds for every non-null $A\subseteq H'$.

\emph{Step 5: a contradiction.} Let $E$ be a two-sided selection on
$(e_1,e_3)$; its thresholds in world~$1$'s frame satisfy
$0<\beta^{(1)}_E\le\alpha^{(1)}_E<\infty$. The set
\[
 T\coloneqq O_1\cap\{t^1_3<\beta^{(1)}_E\}\cap\{z_0+1<t^1_2<z_0+2\}
\]
is non-null by Lemma~\ref{lem:regions}(c), and $T\subseteq R$, since
$O_1\subseteq D_{12}$. We apply Theorem~\ref{thm:own} with
$(u,k,l)\coloneqq(3,1,2)$, three distinct worlds of $S(K)$, and the
two-sided selection $E$ on $(e_3,e_1)$. In world~$3$'s frame $E$ has
thresholds $(\alpha^{(3)}_E,\beta^{(3)}_E)=(1/\beta^{(1)}_E,
1/\alpha^{(1)}_E)$, so the set $T_0$ of Theorem~\ref{thm:own} is
$O_1\cap\{t^3_1>1/\beta^{(1)}_E\}=O_1\cap\{t^1_3<\beta^{(1)}_E\}\supseteq
T$, and its sets $Y_c$ are $O_2\cap\Hi21c=O_2\cap\{t^1_2<1/c\}$. So there
is $c\in(0,\infty)$ such that $T$ is not a finite union of measurable sets
$Q$ for which $(Q,\Gamma)$ is a weakly SI pair of world~$1$, where
$\Gamma\coloneqq O_2\cap\{t^1_2<1/c\}$. The set $\Gamma$ is non-null
(Lemma~\ref{lem:regions}(c)), so
Theorem~\ref{thm:LH}(i) gives an SI pair $(A,\Gamma')$ of world~$1$ with $A$
non-null and $\Gamma'\subseteq\Gamma$. By Step~4, $R\in\Jid_A$; so, up to a
null set, $T$ is contained in a finite union $Q_1\cup\dots\cup Q_r$ of
measurable sets $Q_i\preceq_dA$. Put $Q'_i\coloneqq Q_i\cap T$, and add the
null set $T\setminus(Q'_1\cup\dots\cup Q'_r)$ to $Q'_1$, which changes no
model. Then $T=Q'_1\cup\dots\cup Q'_r$, and $Q'_i\preceq_dQ_i\preceq_dA$ by
weak (G). Each $Q'_i\subseteq O_1$ misses $\Gamma\subseteq O_2$, and $A$
misses $\Gamma'$; so the model $V+\cred_1(\Gamma')$ is silent on
$Q'_i\cup A$, and
\[
 V+\debt_1(Q'_i)+\cred_1(\Gamma)\preceq_1V+\debt_1(Q'_i)+\cred_1(\Gamma')
 \preceq_1V+\debt_1(A)+\cred_1(\Gamma')\prec_1V.
\]
The first step is weak (G) (more credit at $1$ on
$\Gamma\setminus\Gamma'$), the second is $Q'_i\preceq_dA$ in the context
$V+\cred_1(\Gamma')$, by (IB$'$), and the third says that $(A,\Gamma')$ is
SI. So every pair $(Q'_i,\Gamma)$ is SI, in particular weakly SI, and $T$
is their union. This contradicts the choice of $c$.
\end{proof}

Steps~1--4 use (G) at every world, completeness and (IB) at world~$1$, and
(S), but nothing about $W_0$: if every selection $C(s)$ with $s<s_e$ near
$s_e$ has $e_2$ in its credal set, then world~$1$'s debt order is fine on a
band region of $D_{12}$. Step~5 excludes this when $W_0$ is undominated;
it is the only step that uses Theorems~\ref{thm:LH} and~\ref{thm:own} and
the two-sided selection on $[e_1,e_3]$. Steps~1--4 are those of
the three-world note. In its Step~5 the set $T$ lies in $P_1$, the credit
set $\Gamma$ in $P_2$, and the inputs are FSE's Corollary~31 and
Theorems~56 and~64; here $T\subseteq O_1$, $\Gamma\subseteq O_2$, the
two-sided selection is a hypothesis, and the inputs are
Theorems~\ref{thm:LH} and~\ref{thm:own}.

\section{Light sets and counting}\label{sec:counting}

In this section we work at world~$2$ and on two of its edges: on $D_{21}$
with the ratio $t^2_1=1/t^1_2$, and on $D_{23}$ with the ratio $t^2_3$.
$O_2\subseteq D_{21}\cap D_{23}$, and by Lemma~\ref{lem:regions}(c) the
joint law of $(t^2_1,t^2_3)$ on $O_2$ has a positive density on
$(0,\infty)^2$. Definition~\ref{def:light} and
Lemmas~\ref{lem:compose} and~\ref{lem:buys} are the three-world note's
Definition~7 and Lemmas~9 and~10. Their statements and proofs concern
world~$2$'s ranking only and do not use $n=3$; we repeat the proofs, since
that note states them under its standing assumption $n=3$.

\begin{definition}\label{def:light}
For measurable $W,T\subseteq H'$, the set $W$ is \emph{$T$-light} if
$W\cap T=\varnothing$ and $D_2+\debt_2(W)\preceq_2D_2+\sil(T)$.
\end{definition}

So the full change from credit to debt at world~$2$ on $W$ costs world~$2$
no more than the loss of its credit on $T$. By weak (G), every measurable
subset of a $T$-light set is $T$-light.

\begin{lemma}\label{lem:near}
Assume (G) at worlds $1$ and $2$ and (S), and suppose that the far end of
the frame $(1,2)$ is rich. Let $c>0$ and let $T\subseteq D_{21}\cap
\{t^2_1>c\}$ be non-null. Then there is $b>0$ such that every measurable
$W\subseteq D_{21}\cap\{t^2_1<b\}$ is $T$-light.
\end{lemma}
\begin{proof}
This is the three-world note's Lemma~8, with richness as a hypothesis; its
proof uses only the edge $[e_1,e_2]$. Let $q_e$ be the far end of the frame
$(1,2)$. By richness there are $q_m$ in the open edge $(e_1,e_2)$ with
$q_m\to q_e$ and $\alpha^{(1)}_{q_m}<\infty$. In world~$2$'s frame on
$[e_2,e_1]$, $C_{q_m}$ has thresholds
$\alpha_m\coloneqq\alpha^{(2)}_{q_m}=1/\beta^{(1)}_{q_m}$ and
$\beta_m\coloneqq\beta^{(2)}_{q_m}=1/\alpha^{(1)}_{q_m}>0$. Since
$C_{q_e}=D_2$, \ref{in:ends}(a) gives $\beta^{(1)}_{q_m}\to\infty$, so
$\alpha_m\to0$.

Let $M_T\coloneqq D_2+\sil(T)$. On $T\subseteq D_{21}$, $D_2$ gives credit
at $2$, which by \eqref{eq:swap} is debt at $1$, while $M_T$ is silent;
elsewhere the two models agree. So $M_T\prec_1D_2=C_{q_e}$ by strict (G)
at world~$1$, and (T2) with $(v,w)=(2,1)$ gives $C_q\preceq_2M_T$ for all $q$
in the open edge near $q_e$. Fix $m$ with $C_{q_m}\preceq_2M_T$ and
$\alpha_m<c$, and put $b\coloneqq\beta_m$; then $0<b\le\alpha_m<c$.

Let $W\subseteq D_{21}\cap\{t^2_1<b\}$ be measurable. The selection
$C_{q_m}$ gives debt at $2$ on $D_{21}\cap\{t^2_1<b\}\supseteq W$, and
$D_2+\debt_2(W)$ has world-$2$ rank $2$ on $W$ and $0$ elsewhere on $H'$.
So $D_2+\debt_2(W)\le^2_\pw C_{q_m}$, and weak (G) gives
$D_2+\debt_2(W)\preceq_2C_{q_m}\preceq_2M_T$. Finally $W\cap T=\varnothing$,
since $W\subseteq\{t^2_1<b\}$, $T\subseteq\{t^2_1>c\}$ and $b<c$.
\end{proof}

\begin{lemma}\label{lem:compose}
Assume (IB) at world~$2$. Let $W_1,\dots,W_N$ be pairwise disjoint
measurable sets, and $T_1,\dots,T_N$ pairwise disjoint measurable sets with
$T_i\cap W_{i'}=\varnothing$ for all $i,i'$. If each $W_i$ is $T_i$-light,
then $W_1\cup\dots\cup W_N$ is $(T_1\cup\dots\cup T_N)$-light.
\end{lemma}
\begin{proof}
For $0\le i\le N$ put
$M_i\coloneqq D_2+\sil(T_1\cup\dots\cup T_i)+\debt_2(W_{i+1}\cup\dots\cup
W_N)$; all the sets involved are pairwise disjoint, so $M_i$ is well
defined, $M_0=D_2+\debt_2(W_1\cup\dots\cup W_N)$ and
$M_N=D_2+\sil(T_1\cup\dots\cup T_N)$. For $0\le i<N$, the models $M_i$ and
$M_{i+1}$ differ exactly on $W_{i+1}\cup T_{i+1}$, with world-$2$ ranks $2$
and $0$ on $W_{i+1}$, and $0$ and $1$ on $T_{i+1}$. So do
$D_2+\debt_2(W_{i+1})$ and $D_2+\sil(T_{i+1})$, and the first is
$\preceq_2$ the second because $W_{i+1}$ is $T_{i+1}$-light. By (IB$'$),
$M_i\preceq_2M_{i+1}$. By transitivity $M_0\preceq_2M_N$, and
$\bigcup_iW_i$ misses $\bigcup_iT_i$.
\end{proof}

\begin{lemma}\label{lem:buys}
Assume (G) and (IB) at world~$2$. Let $W,T,Y\subseteq H'$ be pairwise
disjoint measurable sets such that $W$ is $T$-light and $T\preceq^2_cY$.
\begin{enumerate}[label=(\alph*),leftmargin=*,itemsep=2pt]
\item $(W,Y)$ is a weakly SI pair of world~$2$.
\item If moreover $P\subseteq H'$ is measurable and disjoint from
$W\cup T\cup Y$, and $(P,Y)$ is a strictly AC pair of world~$2$, then
$D_2+\debt_2(W)\prec_2D_2+\debt_2(P)$.
\end{enumerate}
\end{lemma}
\begin{proof}
(a) By lightness and weak (G),
$D_2+\debt_2(W)\preceq_2D_2+\sil(T)\preceq_2D_2+\sil(T\cup W)$. The outer
two models differ exactly on $W\cup T$, with world-$2$ ranks $2$ and $1$ on
$W$, and $0$ and $1$ on $T$; so do $V+\debt_2(W)+\cred_2(T)$ and $V$. By
(IB$'$), $V+\debt_2(W)+\cred_2(T)\preceq_2V$. By definition,
$T\preceq^2_cY$ means $V+\cred_2(Y)\preceq_2V+\cred_2(T)$. The model
$V+\debt_2(W)$ is silent on $T\cup Y$, so (IB$'$) gives
$V+\debt_2(W)+\cred_2(Y)\preceq_2V+\debt_2(W)+\cred_2(T)\preceq_2V$.

(b) We show
\begin{multline*}
 D_2+\debt_2(W)\preceq_2D_2+\sil(T)\preceq_2D_2+\sil(Y)\\
 \preceq_2D_2+\sil(Y\cup P)\prec_2D_2+\debt_2(P).
\end{multline*}
The first step is lightness. In the second, the two models differ exactly
on $T\cup Y$, with ranks $1$ and $0$ on $T$, and $0$ and $1$ on $Y$; so do
$V+\cred_2(Y)$ and $V+\cred_2(T)$, and (IB$'$) turns $T\preceq^2_cY$ into
the second step. The third step is weak (G). In the fourth, the two models
differ exactly on $Y\cup P$, with ranks $1$ and $0$ on $Y$, and $1$ and $2$
on $P$; so do $V$ and $V+\debt_2(P)+\cred_2(Y)$, and (IB$'$) turns the
strict AC of $(P,Y)$ into the fourth step.
\end{proof}

For $c>0$ put $Y_c\coloneqq O_3\cap\{t^2_3<1/c\}=O_3\cap\Hi32c$, a non-null
set; for $(u,k,l)=(1,2,3)$ this is the set $Y_c$ of Theorem~\ref{thm:own}. Let $\Iid_c$ be the family of measurable sets that are contained, up
to a null set, in a finite union of measurable sets $Q\subseteq O_2$ for
which $(Q,Y_c)$ is a weakly SI pair of world~$2$; such $Q$ miss
$Y_c\subseteq O_3$. $\Iid_c$ is an ideal: it contains the null sets and is
closed under measurable subsets and finite unions. If $c\le c'$, then
$Y_{c'}\subseteq Y_c$, and a weakly SI pair $(Q,Y_{c'})$ gives a weakly SI
pair $(Q,Y_c)$ by weak (G) at world~$2$ (more credit); so
$\Iid_{c'}\subseteq\Iid_c$. Let $\UU_c$ be the set of ultrafilters that
contain $O_2$ and avoid $\Iid_c$.

If world~$2$ is complete, $\Uf\in\UU_c$, $P\in\Uf$ and $P\subseteq O_2$,
then $(P,Y_c)$ is a strictly AC pair of world~$2$. Indeed $P\notin\Iid_c$,
so $(P,Y_c)$ is not weakly SI, since otherwise $P$ would cover itself; and
by completeness a pair that is not weakly SI is strictly AC.

\begin{proposition}\label{prop:finite}
Let $n\ge3$, assume (G) at every world and (S), and let world~$2$ be
complete with (IB). Suppose that the far end of the frame $(1,2)$ is rich.
Then $\UU_c$ is finite for every $c>0$.
\end{proposition}
\begin{proof}
Fix $c>0$.

\emph{Credit witnesses.} By \ref{in:yard}(b) at world~$2$ for the edge
$[e_2,e_1]$, applied to the non-null set $Y_c$, there is $a>0$ such that
every measurable subset of $\Lo21a=D_{21}\cap\{t^2_1<a\}$ is
$\preceq^2_cY_c$. Put $T\coloneqq O_1\cap\{a/2<t^2_1<a\}$; it is non-null by
Lemma~\ref{lem:regions}(c), and every measurable $T'\subseteq T$ satisfies
$T'\preceq^2_cY_c$. Split $T$ into pairwise disjoint non-null measurable
sets $T_j$, $j\ge1$; this is possible since $\mu$ is atomless. Each $T_j$
lies in $D_{21}\cap\{t^2_1>a/2\}$, so Lemma~\ref{lem:near} gives $b_j>0$
such that every measurable subset of $D_{21}\cap\{t^2_1<b_j\}$ is
$T_j$-light.

\emph{A band-dense family of light sets.} Choose reals $\varepsilon_j$ with
$0<\varepsilon_1<b_1$ and $0<\varepsilon_j<\min(b_j,\varepsilon_{j-1}/2)$
for $j\ge2$. Let $(r_j,r'_j)_{j\ge1}$ enumerate the pairs of rationals
with $0<r_j<r'_j$, and put
\[
 W_j\coloneqq O_2\cap\{r_j<t^2_3<r'_j\}\cap
 \{\varepsilon_j/2<t^2_1<\varepsilon_j\}.
\]
Each $W_j$ is non-null by Lemma~\ref{lem:regions}(c). The $W_j$ are
pairwise disjoint, since their ranges of $t^2_1$ are. Each $W_j\subseteq
O_2$ misses $T\subseteq O_1$ and $Y_c\subseteq O_3$, and $W_j$ is
$T_j$-light, since $W_j\subseteq D_{21}\cap\{t^2_1<b_j\}$. Let $\Bfam$ be
the family of measurable subsets of the sets $W_j$. It is band-dense for
the edge $[e_2,e_3]$ in world~$2$'s frame: a band
$D_{23}\cap\{s<t^2_3<s'\}$, with $0\le s<s'\le\infty$, contains every $W_j$
with $s\le r_j<r'_j\le s'$, and there are such $j$.

\emph{The chain.} By \ref{in:cut} with $(v,w)=(2,3)$ (world~$2$ is
complete, (G) holds at worlds $2$ and $3$, and (S) holds), there is a
connecting chain from $D_2$ to $D_3$ with perturbation sets in $\Bfam$. Let
$\ell\ge2$ be its length and $B_1,\dots,B_{\ell-1}$ its perturbation sets.
Each $B_i$ lies in some $W_{j_i}$; let $N\coloneqq\max_ij_i$ and
\[
 F\coloneqq W_1\cup\dots\cup W_N,\qquad T^N\coloneqq T_1\cup\dots\cup T_N.
\]
By Lemma~\ref{lem:compose}, $F$ is $T^N$-light, and so is every
$B_i\subseteq F$. Moreover $T^N\preceq^2_cY_c$, and $F\subseteq O_2$,
$T^N\subseteq O_1$ and $Y_c\subseteq O_3$ are pairwise disjoint. By
Lemma~\ref{lem:buys}(a), $(F,Y_c)$ is a weakly SI pair of world~$2$; as
$F\subseteq O_2$, $F\in\Iid_c$.

\emph{Counting.} Suppose that $\UU_c$ has $\ell$ distinct members
$\Uf^1,\dots,\Uf^\ell$. Each contains $O_2$ and not $F$, hence contains
$O_2\setminus F$. For $m<m'$ choose $G_{mm'}\in\Uf^m$ with
$H'\setminus G_{mm'}\in\Uf^{m'}$, and put $G_{m'm}\coloneqq H'\setminus
G_{mm'}$; then $G_{mm'}\in\Uf^m$ for all $m\ne m'$. The sets
\[
 P_m\coloneqq(O_2\setminus F)\cap\bigcap_{m'\ne m}G_{mm'}\in\Uf^m
\]
are pairwise disjoint: for $m\ne m'$, $P_m\subseteq G_{mm'}$ and
$P_{m'}\subseteq G_{m'm}=H'\setminus G_{mm'}$. Each $P_m$ is a subset of
$O_2$ in $\Uf^m\in\UU_c$, so $(P_m,Y_c)$ is a strictly AC pair of
world~$2$. For all $i$ and $m$, the sets $B_i\subseteq F$,
$T^N\subseteq O_1$, $Y_c\subseteq O_3$ and $P_m\subseteq O_2\setminus F$
are pairwise disjoint, $B_i$ is $T^N$-light, and $T^N\preceq^2_cY_c$. So
Lemma~\ref{lem:buys}(b) gives
\[
 D_2+\debt_2(B_i)\prec_2D_2+\debt_2(P_m)\qquad
 \text{for all }1\le i\le\ell-1\text{ and }1\le m\le\ell.
\]
Since $P_m\subseteq O_2\setminus F\subseteq D_{23}\setminus F$, this
contradicts \ref{in:count} with $(v,w)=(2,3)$. Hence $|\UU_c|\le\ell-1$.
\end{proof}

This is the three-world note's Proposition~11 with its three-world regions
replaced by $O_1,O_2,O_3$. Only world~$2$, the edges $[e_1,e_2]$ and $[e_2,e_3]$, and
the disjointness of the three regions are involved; no dominance argument
is used, and nothing is assumed about $V$ or $W_0$ beyond what richness
requires.

\begin{theorem}\label{thm:count}
Assume (P1$_n$). Let $W_0=D_K$ be a coherent model with $1,2,3\in S(K)$,
and suppose that the edges $[e_1,e_2]$ and $[e_1,e_3]$ carry two-sided
selections. Then $W_0$ is dominated.
\end{theorem}
\begin{proof}
Suppose that $W_0$ is undominated. By Proposition~\ref{prop:rich}, the far
end of the frame $(1,2)$ is rich. By Proposition~\ref{prop:finite}, the set
$\UU\coloneqq\bigcup_{c\in\mathbb N,\,c\ge1}\UU_c$ is countable; list its
members as $\Uf^1,\Uf^2,\dots$ (the list may be finite or empty). Let $E$
be a two-sided selection on $(e_1,e_2)$. As $\alpha^{(1)}_E<\infty$, the
set $T_0\coloneqq O_2\cap\{t^1_2>\alpha^{(1)}_E\}$ is non-null
(Lemma~\ref{lem:regions}(c)). By
\ref{in:uf} choose $B^i\in\Uf^i$ with $\mu(B^i)<2^{-i-1}\mu(T_0)$, and put
$T^*\coloneqq T_0\setminus\bigcup_iB^i$. Then $\mu(T^*)\ge\mu(T_0)/2>0$, and
$T^*\cap B^i=\varnothing$ for every $i$.

Theorem~\ref{thm:own} with $(u,k,l)\coloneqq(1,2,3)$ and the selection $E$
applies: its set $T_0$ is the present $T_0$, and its sets $Y_c=O_3\cap\Hi
32c$ are the present $Y_c$. It gives $c_0$ such that, for every integer
$c\ge c_0$, $T^*$ is not a finite union of measurable sets $Q$ for which
$(Q,Y_c)$ is weakly SI at world~$2$. Fix such an integer $c\ge1$.

Then $T^*\notin\Iid_c$. Otherwise, up to a null set,
$T^*\subseteq Q_1\cup\dots\cup Q_r$ with $Q_i\subseteq O_2$ and $(Q_i,Y_c)$
weakly SI; here $r\ge1$, since $T^*$ is non-null. Then $T^*$ is the union
of the sets $Q_i\cap T^*$ together with a null set, which we add to
$Q_1\cap T^*$; this changes no model. Each $(Q_i\cap T^*,Y_c)$ is weakly SI
by weak (G) at world~$2$ (less debt), contrary to the choice of $c$.

By \ref{in:uf} there is an ultrafilter $\Uf$ that contains $T^*$ and avoids
$\Iid_c$. It contains $O_2\supseteq T^*$, so $\Uf\in\UU_c\subseteq\UU$, and
$\Uf=\Uf^i$ for some $i$. Then $\Uf$ contains the disjoint sets $T^*$ and
$B^i$, which is impossible.
\end{proof}

For $n=3$ and $W_0=V$, with the two-sided selections supplied by FSE's
Corollary~31, this is the three-world note's Theorem~12.

\section{Main results}\label{sec:main}

\begin{theorem}\label{thm:clique}
Let $n\ge3$, let every ranking be a complete preorder satisfying (G) and
(IB), and assume (S). Let $K$ be the credal set of a coherent model $D_K$,
and suppose that $S(K)=\{x:e_x\in K\}$ has at least three members. If
$D_K$ is undominated, then (EP$_{xy}$) holds for all distinct $x,y\in
S(K)$, and every edge $[e_x,e_y]$ with $x,y\in S(K)$ is fine.
\end{theorem}
\begin{proof}
Put $W_0\coloneqq D_K$. By Theorem~\ref{thm:dichotomy}, either (EP$_{xy}$)
holds for all distinct $x,y\in S(K)$ or for none. Suppose it holds for
none, and label three distinct worlds of $S(K)$ as $1,2,3$. The edges
$[e_1,e_2]$ and $[e_1,e_3]$ carry selections whose credal sets omit both
endpoints, and these are two-sided. Theorem~\ref{thm:count} makes $W_0$
dominated, contrary to the hypothesis. So (EP$_{xy}$) holds for all
distinct $x,y\in S(K)$, and \ref{in:EPfine} makes each such edge fine.
\end{proof}

\begin{theorem}\label{thm:main}
Let $n\ge3$, let every ranking be a complete preorder satisfying (G) and
(IB), and assume (S). Then $V$ is dominated, and hence dominated by a
selected coherent model.
\end{theorem}
\begin{proof}
Suppose that $V$ is undominated. Theorem~\ref{thm:clique} with
$K\coloneqq\Delta$, for which $S(K)$ is the set of all worlds, gives
(EP$_{ij}$) for every edge. Then \ref{in:reduction} makes every coherent
model whose credal set contains three vertices dominated, in particular
$V$: a contradiction. So $V$ is dominated, and by CBOI's Lemma~1 it is
dominated by a selected coherent model.
\end{proof}

Equivalently, as in the three-world note's Corollary~13(a): if (G) and (IB)
hold at every world and (S) holds, and $V$ is undominated, which is the
case in particular if every coherent model is undominated, then the
ranking of some world is incomplete.

\begin{corollary}\label{cor:consequences}
Let $n\ge3$, let every ranking be a complete preorder satisfying (G) and
(IB), and assume (S).
\begin{enumerate}[label=(\alph*),leftmargin=*,itemsep=2pt]
\item If some edge $[e_x,e_y]$ carries a selection whose credal set omits
both endpoints, then every coherent model whose credal set contains $e_x$,
$e_y$ and at least one further vertex is dominated.
\item If some coherent model $D_K$ with $|S(K)|\ge3$ is undominated, then
the edges with the endpoint property form a graph on the set of worlds
that is not complete and in which $S(K)$ is a clique. In particular
$S(K)$ is not the set of all worlds, $K\ne\Delta$, and $n\ge4$.
\item If $n=3$, every coherent model whose credal set contains three
vertices is dominated.
\end{enumerate}
\end{corollary}
\begin{proof}
(a) If $D_K$ is such a model and is undominated, then $x,y\in S(K)$ and
$|S(K)|\ge3$, so Theorem~\ref{thm:clique} gives (EP$_{xy}$), which the
given selection violates. (b) By Theorem~\ref{thm:clique}, $S(K)$ is a
clique. If every edge had the endpoint property, \ref{in:reduction} would
make $D_K$ dominated. So some edge fails it, and that edge does not join
two worlds of $S(K)$; hence $S(K)$ is not the set of all worlds, and as
$|S(K)|\ge3$, $n\ge4$. (c) For $n=3$, a convex set containing $e_1,e_2,e_3$
is $\Delta$, so the only such model is $V$, which is dominated by
Theorem~\ref{thm:main}.
\end{proof}

\section{Remarks}\label{sec:remarks}

\begin{remark}\label{rem:hyp}
\emph{Where the hypotheses are used.}
\begin{enumerate}[label=(\alph*),leftmargin=*,itemsep=2pt]
\item \emph{Completeness.} Through (T1) (Lemma~\ref{lem:EPnotCN},
Proposition~\ref{prop:rich}, and the no-saturated-cut theorem \ref{in:cut}
at world~$2$); through \ref{in:yard}(a), \ref{in:SAC}, \ref{in:ext} and
\ref{in:EPfine}; in reading AC pairs as $V\preceq_vV+\debt_v(A)+\cred_v(T)$
(Lemmas~\ref{lem:conestails}(a), \ref{lem:CNtoDI}); and in turning the
failure of $X\preceq Y$ into $Y\prec X$ (Step~1 of Theorem~\ref{thm:own}
and of Proposition~\ref{prop:partner}, and the strict AC of the members of
ultrafilters in $\UU_c$). Theorem~\ref{thm:LH} uses the completeness of
every world. Lemmas~\ref{lem:XOR}, \ref{lem:near}, \ref{lem:compose},
\ref{lem:buys}, \ref{in:yard}(b), (T2) and the counting lemma
\ref{in:count} do not use completeness.
\item \emph{(IB).} At every world, through Theorem~\ref{thm:LH}, and at
the worlds of each triple in Sections~\ref{sec:own}--\ref{sec:counting}.
Lemma~\ref{lem:near} and \ref{in:cut} use no (IB).
\item \emph{(G).} Strict (G) is used at every world: in
Lemma~\ref{lem:regions}(e), that is, at the worlds outside the active
triple of each dominance step, and through the results of FSE.
Lemmas~\ref{lem:compose} and~\ref{lem:buys} and \ref{in:count} need only
weak (G).
\item \emph{Undominatedness.} Section~\ref{sec:LH} assumes nothing about
any model. Undominatedness of $W_0=D_K$ is the reductio hypothesis of
Theorems~\ref{thm:clique} and~\ref{thm:main}. It enters only through the
two dominance steps, Step~1 of Theorem~\ref{thm:own} and Step~1 of
Proposition~\ref{prop:partner}; hence through Theorem~\ref{thm:dichotomy},
Step~5 of Proposition~\ref{prop:rich}, and Theorem~\ref{thm:count}.
\item \emph{(S).} Through (T1), (T2), the face lemma, the competitiveness
of selections on open edges on their own two worlds, and the results of
FSE, CCF and the three-world note. The selections used are all on edges;
higher-dimensional faces enter only inside \ref{in:reduction} and CBOI's
Lemma~1, which are stated for every $n\ge3$.
\end{enumerate}
\end{remark}

\begin{remark}\label{rem:regions}
\emph{Why the regions $O_v$.} There are exactly two dominance
constructions: Step~1 of the proof of Theorem~\ref{thm:own}, with
$d=\omega_k(\Sigma_{i^*})+\omega_u(Z)+\omega_l(Y)$, and Step~1 of the proof
of Proposition~\ref{prop:partner}, with
$d=\omega_v(\Sigma_{i^*})+\omega_u(B)+\omega_w(Z)$. Each changes $W_0$ only
on subsets of $O_x$, for $x$ in a triple of worlds of $S(K)$, by the option
$\omega_x$, which is good at every world other than $x$. So at every world
outside the triple each piece of the change is a change from silence to
credit, and $W_0+d\prec W_0$ there (Lemma~\ref{lem:regions}(e)). For
$n=3$ there are no such worlds and $O_x=P_x$. For $n\ge4$ the three-world
regions would not do: on $D_{ku}\cap D_{kl}\setminus O_k$ some world
$x\notin\{u,k,l\}$ has the sign of $k$, and the option bad at $k$ is bad
at $x$ as well, so a change made of it would not be weakly better than
$W_0$ at $x$. In Section~\ref{sec:LH}, on the other hand, the sets
$D_{wv}\cap D_{wu}$ and $D_{vw}\cap D_{vu}$ are used in comparisons at a
single world and in the competitiveness of an edge selection on its own
two worlds, where any regions with the right signs at the worlds involved
suffice. A further constraint is that the base model must be silent on the
regions it is changed on; for the regions $O_x$, $x\in S(K)$, this is
Lemma~\ref{lem:regions}(f).
\end{remark}

\begin{remark}\label{rem:open}
\emph{What remains for coherent models with three vertices.} The conclusion
of CBOI's Theorem~2 and of CCF's Theorem~16 is that \emph{every} coherent
model whose credal set contains three vertices is dominated, not only $V$.
For $n=3$ the two coincide (Corollary~\ref{cor:consequences}(c)). For
$n\ge4$, under (P1$_n$), Theorem~\ref{thm:clique} and
Corollary~\ref{cor:consequences}(b) reduce the stronger conclusion to the
following case, which is open: $D_K$ undominated, $3\le|S(K)|<n$,
(EP$_{xy}$) on every edge with $x,y\in S(K)$, and failure of (EP) on some
edge leaving $S(K)$ or joining two worlds outside it. The arguments of
Sections~\ref{sec:own}--\ref{sec:counting} do not reach the edges that
leave $S(K)$, for the following reason. Let $x\notin S(K)$, let $g\in O_x$,
and let $h\in\{g,-g\}$ represent the option $\omega_x$, so that $h_x<0$ and
$h_y>0$ for every $y\ne x$. For $s\in S(K)$, $e_s\cdot(-h)=-h_s<0$, so
$D_K$ never recommends $-h$: on $O_x$ it gives either silence or
$\omega_x$. Since $K$ is compact and $e_x\notin K$, $\rho\coloneqq
\max_{p\in K}p_x<1$, and $p\cdot h\ge-\rho|g_x|+(1-\rho)\min_{y\ne x}|g_y|$
for $p\in K$; so $D_K$ gives $\omega_x$ on the nonempty open, hence
non-null, set of $g\in O_x$ with $\rho|g_x|<(1-\rho)\min_{y\ne x}|g_y|$. If
$K$ lies in the face spanned by $\{e_s:s\in S(K)\}$, then $\rho=0$, and
$D_K$ gives $\omega_x$ on all of $O_x$. Now the dominance step of
Proposition~\ref{prop:partner} for an edge $[e_x,e_v]$, $v\in S(K)$, adds
$\omega_x$ on a set $B\subseteq O_x$ and appeals to Lemma~\ref{lem:XOR}
with $M\coloneqq W_0$, which requires $W_0$ to be silent on $B$. Where
$W_0$ already gives $\omega_x$, the part of the change on $B$ is not a
change at all: the pairs $(E,E')$ and $(W_0+\tau_{T,B},W_0)$ of
Lemma~\ref{lem:XOR} no longer have the same difference set, and the
competitiveness of $E$ says nothing about $W_0$. If $K$ lies in the face
spanned by $\{e_s:s\in S(K)\}$, this happens on all of $O_x$, so the step
cannot start. The same applies to Theorem~\ref{thm:own} for triples that
are not contained in $S(K)$. And
CCF's Theorem~16, which closes the case $K=\Delta$, needs the endpoint
property on every edge: through CCF's Theorem~15 it gives every world a
weak extension of CBOI's form, and CBOI's construction (CCF, Lemma~1)
needs such an extension at every world. With (EP) only on the edges
between worlds of $S(K)$, the worlds of $S(K)$ are fine on the regions
$D_{xy}$, $x,y\in S(K)$ (\ref{in:EPfine}), but nothing is known about their
debt and credit orders on $D_{xz}$ with $z\notin S(K)$, nor about the
worlds outside $S(K)$. By (S), every selection $C_p$ is coherent and
undominated, so Theorem~\ref{thm:clique} and
Corollary~\ref{cor:consequences}(b) apply to it: if the credal set of $C_p$
contains at least three vertices, the endpoint property holds on all edges
between them and fails on some other edge, and $n\ge4$. (For $n=3$, no
selection has three vertices in its credal set.) So the open case contains
the concrete question whether a selection $C_p$, with $p$ in the relative
interior of a face with at least three vertices, can have three of these
vertices in its credal set. What is proved is the clique constraint: an
undominated coherent model with three vertices in its credal set can exist
only if the endpoint property holds on all edges between those vertices
and fails on some other edge; and one edge with a selection that omits both
endpoints makes every coherent model whose credal set contains that edge's
endpoints and a third vertex dominated.
\end{remark}

\end{document}
